% =====================================================================
%  شرح الشاطبية — حرز الأماني ووجه التهاني
%  الملف الرئيسي. يُترجَم بـ XeLaTeX.
%
%  لإضافة باب جديد:
%    ١) انسخ chapters/_template.tex باسم جديد (مثلًا chapters/05-idgham-kabir.tex)
%    ٢) أضف سطر \include{chapters/05-idgham-kabir} في المكان المناسب أدناه
%  لترجمة باب واحد فقط أثناء الكتابة (أسرع): فعّل \includeonly أدناه.
% =====================================================================
\documentclass[12pt,oneside]{book}

% =====================================================================
%  setup.tex — الحزم والخطوط وتخطيط الصفحة
%  يُترجَم بـ XeLaTeX فقط (في Overleaf: Menu > Compiler > XeLaTeX)
% =====================================================================

% ---------- تخطيط الصفحة ----------
\usepackage[a4paper,top=2.4cm,bottom=2.4cm,inner=2.2cm,outer=2.2cm,headheight=16pt]{geometry}

% ---------- حزم عامة (تُحمَّل قبل polyglossia/bidi) ----------
\usepackage{fontspec}
\usepackage{xcolor}
\usepackage{graphicx}
\usepackage{array}
\usepackage{tabularx}
\usepackage{booktabs}
\usepackage{multirow}
\usepackage{longtable}
\usepackage{enumitem}
\usepackage{fancyhdr}
\usepackage{titlesec}
\usepackage{titletoc}
\usepackage[most]{tcolorbox}
\usepackage{etoolbox}
\usepackage{needspace}
\usepackage[hidelinks,bookmarksnumbered,unicode]{hyperref}

% ---------- اللغة: العربية لغة رئيسية ----------
\usepackage{polyglossia}
\setmainlanguage[numerals=mashriq]{arabic}
\setotherlanguage{english}

% ---------- الخطوط (مضمّنة في مجلد fonts/ لتعمل على أي جهاز وعلى Overleaf) ----------
\defaultfontfeatures{Ligatures=TeX}
\setmainfont{Amiri}[
  Path=fonts/,
  Extension=.ttf,
  UprightFont=*-Regular,
  BoldFont=*-Bold,
  ItalicFont=*-Italic,
  BoldItalicFont=*-BoldItalic,
  Script=Arabic,
  Scale=1.0]
\newfontfamily\arabicfont{Amiri}[
  Path=fonts/,
  Extension=.ttf,
  UprightFont=*-Regular,
  BoldFont=*-Bold,
  ItalicFont=*-Italic,
  BoldItalicFont=*-BoldItalic,
  Script=Arabic]
\newfontfamily\quranfont{AmiriQuran}[Path=fonts/,Extension=.ttf,Script=Arabic,Scale=1.05]
\newfontfamily\englishfont{Amiri}[Path=fonts/,Extension=.ttf,UprightFont=*-Regular,BoldFont=*-Bold,ItalicFont=*-Italic]

% ---------- المسافات ----------
\linespread{1.25}
\setlength{\parindent}{0pt}
\setlength{\parskip}{0.45em}
\setlist{itemsep=2pt,topsep=3pt,leftmargin=*}

% ---------- رأس الصفحة وتذييلها ----------
\pagestyle{fancy}
\fancyhf{}
\fancyhead[C]{\small\color{headcolor}\leftmark}
\fancyfoot[C]{\color{headcolor}\thepage}
\renewcommand{\headrulewidth}{0.4pt}
\renewcommand{\headrule}{\color{gold}\hrule height 0.6pt}
\renewcommand{\chaptermark}[1]{\markboth{#1}{}}
\fancypagestyle{plain}{\fancyhf{}\fancyfoot[C]{\color{headcolor}\thepage}\renewcommand{\headrulewidth}{0pt}}

% ---------- العناوين ----------
% الباب (chapter) بلا رقم ظاهر، بإطار مزخرف
\titleformat{\chapter}[block]
  {\centering\normalfont}
  {}{0pt}
  {\begin{tcolorbox}[enhanced,colback=parchment,colframe=gold,boxrule=1.2pt,
     arc=3mm,width=0.8\linewidth,center,halign=center,
     borderline={0.4pt}{3pt}{gold}]
     \Huge\bfseries\color{maroon}}
  [\end{tcolorbox}]
\titlespacing*{\chapter}{0pt}{-10pt}{18pt}

\titleformat{\section}[block]
  {\needspace{8\baselineskip}\normalfont\Large\bfseries\color{maroon}}
  {}{0pt}
  {}
  [{\color{gold}\titlerule[0.8pt]}]
\titlespacing*{\section}{0pt}{14pt}{8pt}

\titleformat{\subsection}[block]
  {\normalfont\large\bfseries\color{headcolor}}
  {}{0pt}{}
\titlespacing*{\subsection}{0pt}{10pt}{4pt}

\setcounter{secnumdepth}{-1}   % بلا ترقيم للعناوين
\setcounter{tocdepth}{1}       % الفهرس: الأبواب والأقسام

% تنسيق الفهرس
\titlecontents{chapter}[0pt]{\addvspace{6pt}\large\bfseries\color{maroon}}
  {}{}{\titlerule*[0.6pc]{.}\contentspage}
\titlecontents{section}[1.5em]{\normalsize}
  {}{}{\titlerule*[0.6pc]{.}\contentspage}

% =====================================================================
%  colors.tex — الألوان ومعانيها
%  كل قارئ له لون ثابت في الملف كله، وراوياه يأخذان لونه نفسه.
%  لتغيير لون قارئ غيّر القيمة هنا فقط فيتغيّر في كل المستند.
% =====================================================================

% ---------- ألوان القراء السبعة ----------
\definecolor{cNafi}{HTML}{1565C0}     % نافع وراوياه (قالون، ورش)      — أزرق
\definecolor{cIbnKathir}{HTML}{2E7D32}% ابن كثير وراوياه (البزي، قنبل) — أخضر
\definecolor{cAbuAmr}{HTML}{D35400}   % أبو عمرو وراوياه (الدوري، السوسي) — برتقالي
\definecolor{cIbnAmir}{HTML}{6A1B9A}  % ابن عامر وراوياه (هشام، ابن ذكوان) — بنفسجي
\definecolor{cAsim}{HTML}{C62828}     % عاصم وراوياه (شعبة، حفص)        — أحمر
\definecolor{cHamza}{HTML}{00838F}    % حمزة وراوياه (خلف، خلاد)        — فيروزي
\definecolor{cKisai}{HTML}{AD1457}    % الكسائي وراوياه (أبو الحارث، الدوري) — وردي غامق

% ---------- رموز الجماعة (ث خ ذ ظ غ ش، صحبة، عم ...) ----------
\definecolor{cGroup}{HTML}{5D4037}    % بني

% ---------- عناصر الشرح ----------
\definecolor{gold}{HTML}{B8860B}       % الإطارات والزخارف
\definecolor{parchment}{HTML}{FBF6E9}  % خلفية صندوق البيت
\definecolor{maroon}{HTML}{7B1F1F}     % العناوين
\definecolor{headcolor}{HTML}{4E342E}  % رأس الصفحة والعناوين الفرعية
\definecolor{wordbg}{HTML}{FFF0A8}     % تظليل الكلمة القرآنية المختلَف فيها داخل البيت
\definecolor{ayahbracket}{HTML}{B8860B}% أقواس الآيات ﴿ ﴾
\definecolor{vocabframe}{HTML}{546E7A} % صندوق معاني الكلمات
\definecolor{vocabback}{HTML}{F1F5F7}
\definecolor{ruleframe}{HTML}{2E7D32}  % صندوق القاعدة
\definecolor{ruleback}{HTML}{EEF7EE}
\definecolor{exframe}{HTML}{1565C0}    % صندوق الأمثلة
\definecolor{exback}{HTML}{EEF4FB}
\definecolor{tahrirframe}{HTML}{6A1B9A}% صندوق التحريرات (إتحاف البرية)
\definecolor{tahrirback}{HTML}{F5EEF8}
\definecolor{noteframe}{HTML}{E65100}  % صندوق التنبيه
\definecolor{noteback}{HTML}{FFF4E5}
\definecolor{sumframe}{HTML}{7B1F1F}   % صندوق الخلاصة
\definecolor{sumback}{HTML}{FDF2F2}

% =====================================================================
%  qurra.tex — أوامر أسماء القراء والرواة والرموز
%  الاستعمال في النص:  \Asim{} و\Kisai{}  ← يظهر الاسم بلون صاحبه
% =====================================================================

% اسم ملوَّن عام:  \qname{لون}{الاسم}
\newcommand{\qname}[2]{\textcolor{c#1}{\textbf{#2}}}

% ---------- نافع وراوياه ----------
\newcommand{\Nafi}{\qname{Nafi}{نافع}}
\newcommand{\Qalun}{\qname{Nafi}{قالون}}
\newcommand{\Warsh}{\qname{Nafi}{ورش}}
% ---------- ابن كثير وراوياه ----------
\newcommand{\IbnKathir}{\qname{IbnKathir}{ابن كثير}}
\newcommand{\Bazzi}{\qname{IbnKathir}{البزي}}
\newcommand{\Qunbul}{\qname{IbnKathir}{قنبل}}
% ---------- أبو عمرو وراوياه ----------
\newcommand{\AbuAmr}{\qname{AbuAmr}{أبو عمرو}}
\newcommand{\AbiAmr}{\qname{AbuAmr}{أبي عمرو}}
\newcommand{\DuriAbuAmr}{\qname{AbuAmr}{الدوري عن أبي عمرو}}
\newcommand{\Susi}{\qname{AbuAmr}{السوسي}}
% ---------- ابن عامر وراوياه ----------
\newcommand{\IbnAmir}{\qname{IbnAmir}{ابن عامر}}
\newcommand{\Hisham}{\qname{IbnAmir}{هشام}}
\newcommand{\IbnDhakwan}{\qname{IbnAmir}{ابن ذكوان}}
% ---------- عاصم وراوياه ----------
\newcommand{\Asim}{\qname{Asim}{عاصم}}
\newcommand{\Shuba}{\qname{Asim}{شعبة}}
\newcommand{\Hafs}{\qname{Asim}{حفص}}
% ---------- حمزة وراوياه ----------
\newcommand{\Hamza}{\qname{Hamza}{حمزة}}
\newcommand{\Khalaf}{\qname{Hamza}{خلف}}
\newcommand{\Khallad}{\qname{Hamza}{خلاد}}
% ---------- الكسائي وراوياه ----------
\newcommand{\Kisai}{\qname{Kisai}{الكسائي}}
\newcommand{\AbulHarith}{\qname{Kisai}{أبو الحارث}}
\newcommand{\DuriKisai}{\qname{Kisai}{الدوري عن الكسائي}}

% ---------- الباقون ----------
\newcommand{\Baqun}{\textbf{\textcolor{black!70}{الباقون}}}

% ---------- الرموز داخل البيت ----------
% كلمة فيها رمز قارئ:  \rz{Kisai}{راويه}
\newcommand{\rz}[2]{\textcolor{c#1}{\textbf{#2}}}
% كلمة فيها رمز جماعة (ث خ ذ ظ غ ش أو صحبة، عم ...):  \rzg{صحبة}
\newcommand{\rzg}[1]{\textcolor{cGroup}{\textbf{#1}}}
% الكلمة القرآنية المختلف فيها داخل البيت:  \qw{مالك}
\newcommand{\qw}[1]{{\setlength{\fboxsep}{1.2pt}\colorbox{wordbg}{#1}}}

% ---------- الآيات ----------
% \ayah{مَالِكِ يَوْمِ الدِّينِ}  ← ﴿ ... ﴾ بخط المصحف
\newcommand{\ayah}[1]{{\quranfont\textcolor{ayahbracket}{﴿}#1\textcolor{ayahbracket}{﴾}}}
% \src{الفاتحة: ٤}  ← مرجع الآية
\newcommand{\src}[1]{{\small\textcolor{black!60}{[#1]}}}

% حرف رمز داخل الجداول بخلفية ملوّنة:  \rzl{Nafi}{أ}
\newcommand{\rzl}[2]{{\setlength{\fboxsep}{2pt}\colorbox{c#1!15}{\textcolor{c#1}{\textbf{\large #2}}}}}
\newcommand{\rzlg}[1]{{\setlength{\fboxsep}{2pt}\colorbox{cGroup!15}{\textcolor{cGroup}{\textbf{\large #1}}}}}

% سهم «يؤدي إلى» باتجاه القراءة العربية
\newcommand{\ila}{\ $\leftarrow$\ }

% =====================================================================
%  boxes.tex — صناديق الشرح وأمر البيت
% =====================================================================
\tcbset{
  shatibox/.style={enhanced,breakable,arc=2mm,boxrule=0.6pt,
    left=6pt,right=6pt,top=4pt,bottom=4pt,
    fonttitle=\bfseries,coltitle=white,
    before skip=6pt,after skip=6pt},
}

% ---------- البيت ----------
% \bayt{رقم البيت}{الشطر الأول}{الشطر الثاني}
\newcounter{baytnum}
\newcommand{\bayt}[3]{%
  \begin{tcolorbox}[enhanced,colback=parchment,colframe=gold,boxrule=0.8pt,arc=2mm,
      left=4pt,right=4pt,top=5pt,bottom=5pt,before skip=10pt,after skip=6pt,
      borderline west={2.5pt}{0pt}{gold},borderline east={2.5pt}{0pt}{gold}]
    \large
    \begin{tabularx}{\linewidth}{@{}>{\centering\arraybackslash}m{2.2em}
                                  >{\centering\arraybackslash}X
                                  >{\centering\arraybackslash}X@{}}
      {\normalsize\setcounter{baytnum}{#1}\textcolor{gold}{\textbf{(\thebaytnum)}}} & #2 & #3
    \end{tabularx}
  \end{tcolorbox}\label{bayt:#1}}

% بيت طويل يُكتب شطراه سطرين (عند الحاجة)
\newcommand{\baytlong}[3]{%
  \begin{tcolorbox}[enhanced,colback=parchment,colframe=gold,boxrule=0.8pt,arc=2mm,
      left=4pt,right=4pt,top=5pt,bottom=5pt,before skip=10pt,after skip=6pt,
      borderline west={2.5pt}{0pt}{gold},borderline east={2.5pt}{0pt}{gold}]
    \large
    {\normalsize\setcounter{baytnum}{#1}\textcolor{gold}{\textbf{(\thebaytnum)}}}\quad #2\par
    \hfill #3\par
  \end{tcolorbox}\label{bayt:#1}}

% ---------- معاني الكلمات ----------
\newtcolorbox{mufradat}{shatibox,colback=vocabback,colframe=vocabframe,
  title={معاني الكلمات}}
% \kalima{الكلمة}{معناها}
\newcommand{\kalima}[2]{\textbf{\textcolor{vocabframe}{#1}}: #2\par}

% ---------- المعنى ----------
% \maana  أو  \maana[(الأبيات ١١٣–١١٥)]
\newcommand{\maana}[1][]{\par\smallskip\textbf{\textcolor{maroon}{المعنى\ifx\relax#1\relax\else\ #1\fi: }}}

% ---------- القراءات ----------
\newtcolorbox{qiraat}{shatibox,colback=white,colframe=gold!80!black,
  title={\textcolor{white}{القراءات في البيت}},colbacktitle=gold!80!black}
% \qira{القراء}{القراءة}
\newcommand{\qira}[2]{\noindent\textcolor{gold}{$\triangleleft$}\ #1 \;:\; #2\par}

% ---------- القاعدة ----------
\newtcolorbox{qaida}[1][القاعدة]{shatibox,colback=ruleback,colframe=ruleframe,title={#1}}

% ---------- الأمثلة ----------
\newtcolorbox{amthila}[1][أمثلة من القرآن]{shatibox,colback=exback,colframe=exframe,title={#1}}

% ---------- التحريرات ----------
\newtcolorbox{tahrir}[1][تحريرات (من إتحاف البرية بتحريرات الشاطبية)]{shatibox,
  colback=tahrirback,colframe=tahrirframe,title={#1}}
% بيت من الإتحاف داخل صندوق التحريرات
\newcounter{ithafnum}
\newcommand{\ithaf}[3]{\par{\centering\setcounter{ithafnum}{#1}\textcolor{tahrirframe}{\textbf{(\theithafnum)}}\ #2 \quad\ \ #3\par}\smallskip}

% ---------- تنبيه ----------
\newtcolorbox{tanbih}[1][تنبيه]{shatibox,colback=noteback,colframe=noteframe,title={#1}}

% ---------- الخلاصة ----------
\newtcolorbox{khulasa}[1][خلاصة]{shatibox,colback=sumback,colframe=sumframe,title={#1}}

% فاصل بين الأبيات
\newcommand{\fasl}{\par\smallskip{\centering\textcolor{gold}{\large ۞}\par}\smallskip}


% \includeonly{chapters/04-fatiha}   % ← أزل علامة % لترجمة باب واحد فقط

\begin{document}

\frontmatter
% صفحة العنوان
\begin{titlepage}
\centering
\vspace*{1.5cm}
\begin{tcolorbox}[enhanced,colback=parchment,colframe=gold,boxrule=1.5pt,arc=4mm,
  borderline={0.5pt}{4pt}{gold},width=0.9\linewidth,center,halign=center,top=18pt,bottom=18pt]
  {\Large\color{headcolor} الشرح الميسَّر لـ}\par\medskip
  {\fontsize{34}{48}\selectfont\bfseries\color{maroon} حِرْزِ الأَمَانِي وَوَجْهِ التَّهَانِي}\par\medskip
  {\LARGE\color{headcolor} في القراءات السبع}\par\medskip
  {\Large (متن الشاطبية)}\par\bigskip
  {\large للإمام القاسم بن فِيرُّه الشاطبي (ت ٥٩٠ هـ)}
\end{tcolorbox}

\vspace{1.2cm}
{\large شرحٌ ميسَّر بيتًا بيتًا، مع بيان القراء بالألوان، والأمثلة القرآنية، والتحريرات}\par
\vspace{0.8cm}
{\large\color{maroon}\textbf{الجزء الأول}}\par
{\large مقدمة الناظم — باب الاستعاذة — باب البسملة — سورة أم القرآن}\par
{\large (الأبيات ١ – ١١٥)}\par

\vfill
{\normalsize مستخلَص من: الوافي في شرح الشاطبية لعبد الفتاح القاضي، وإتحاف البرية بتحريرات الشاطبية لحسن خلف الحسيني،\\ وتقريب الشاطبية، مع متن الشاطبية المضبوط}\par
\vspace{0.6cm}
\end{titlepage}

\tableofcontents
% =====================================================================
%  دليل الاستخدام: الألوان والرموز والصناديق
% =====================================================================
\chapter{دليل الاستخدام}

\section{طريقة الشرح}
نمشي على الشاطبية \textbf{بيتًا بيتًا}. لكل بيت:
\begin{enumerate}
\item \textbf{نص البيت} في إطار ذهبي، ورقمه بجانبه (بحسب ترقيم المتن المضبوط).
\item \textbf{معاني الكلمات} الصعبة فقط.
\item \textbf{المعنى}: شرح مختصر سهل.
\item في أبواب القراءات: \textbf{صندوق القراءات} يبين من قرأ بكذا، ثم \textbf{القاعدة} إن كانت تحتاج تفسيرًا، ثم \textbf{أمثلة من القرآن}، ثم \textbf{التحريرات} إن وُجدت.
\end{enumerate}
وفي آخر كل باب \textbf{خلاصة} تجمع أحكامه في جدول.

\section{ألوان القراء}
كل قارئ له لون ثابت في المستند كله، وراوياه يأخذان لونه نفسه. فإذا رأيت اللون الأحمر فهو \Asim{} أو أحد راوييه، وهكذا.

\begin{center}
\renewcommand{\arraystretch}{1.5}
\begin{tabular}{c c c c c}
\toprule
\textbf{القارئ} & \textbf{رمزه} & \textbf{الراوي الأول} & \textbf{الراوي الثاني} & \textbf{بلده}\\
\midrule
\Nafi & \rzl{Nafi}{أ} & \Qalun{} \rzl{Nafi}{ب} & \Warsh{} \rzl{Nafi}{ج} & المدينة\\
\IbnKathir & \rzl{IbnKathir}{د} & \Bazzi{} \rzl{IbnKathir}{هـ} & \Qunbul{} \rzl{IbnKathir}{ز} & مكة\\
\AbuAmr & \rzl{AbuAmr}{ح} & \qname{AbuAmr}{الدوري} \rzl{AbuAmr}{ط} & \Susi{} \rzl{AbuAmr}{ي} & البصرة\\
\IbnAmir & \rzl{IbnAmir}{ك} & \Hisham{} \rzl{IbnAmir}{ل} & \IbnDhakwan{} \rzl{IbnAmir}{م} & الشام\\
\Asim & \rzl{Asim}{ن} & \Shuba{} \rzl{Asim}{ص} & \Hafs{} \rzl{Asim}{ع} & الكوفة\\
\Hamza & \rzl{Hamza}{ف} & \Khalaf{} \rzl{Hamza}{ض} & \Khallad{} \rzl{Hamza}{ق} & الكوفة\\
\Kisai & \rzl{Kisai}{ر} & \AbulHarith{} \rzl{Kisai}{س} & \qname{Kisai}{الدوري} \rzl{Kisai}{ت} & الكوفة\\
\bottomrule
\end{tabular}
\end{center}

\begin{tanbih}[طريقة حفظ الرموز]
الحروف مرتبة على «أَبْجَدْ هَوَّزْ حُطِّي كَلَمُنْ سَعْفَصْ قَرَشَتْ»، كل ثلاثة أحرف لقارئ وراوييه بالترتيب:
\textbf{\rzl{Nafi}{أبج}} \textbf{\rzl{IbnKathir}{دهز}} \textbf{\rzl{AbuAmr}{حطي}} \textbf{\rzl{IbnAmir}{كلم}} \textbf{\rzl{Asim}{نصع}} \textbf{\rzl{Hamza}{فضق}} \textbf{\rzl{Kisai}{رست}}\\
(الحرف الأول للإمام، والثاني لراويه الأول، والثالث لراويه الثاني.)
\end{tanbih}

\section{رموز الجماعة}
وهذه رموز لأكثر من قارئ، ولونها \textcolor{cGroup}{\textbf{بنيّ}}:

\begin{center}
\renewcommand{\arraystretch}{1.45}
\begin{tabular}{c l}
\toprule
\textbf{الرمز} & \textbf{القراء} \\
\midrule
\rzlg{ث} & الكوفيون: \Asim، \Hamza، \Kisai \\
\rzlg{خ} & السبعة إلا \Nafi \\
\rzlg{ذ} & \IbnAmir{} والكوفيون \\
\rzlg{ظ} & \IbnKathir{} والكوفيون \\
\rzlg{غ} & \AbuAmr{} والكوفيون \\
\rzlg{ش} & \Hamza{} و\Kisai \\
\midrule
\rzlg{صُحْبَة} & \Hamza، \Kisai، \Shuba \\
\rzlg{صِحَاب} & \Hamza، \Kisai، \Hafs \\
\rzlg{عَمّ} & \Nafi، \IbnAmir \\
\rzlg{سَمَا} & \Nafi، \IbnKathir، \AbuAmr \\
\rzlg{حَقّ} & \IbnKathir، \AbuAmr \\
\rzlg{نَفَر} & \IbnKathir، \AbuAmr، \IbnAmir \\
\rzlg{حِرْمِيّ} & \Nafi، \IbnKathir \\
\rzlg{حِصْن} & \Nafi{} والكوفيون \\
\bottomrule
\end{tabular}
\end{center}

\section{معاني الصناديق والألوان الأخرى}
\begin{center}
\renewcommand{\arraystretch}{1.5}
\begin{tabular}{c l}
\toprule
\textbf{الشكل} & \textbf{معناه} \\
\midrule
\colorbox{parchment}{\textcolor{gold}{\textbf{إطار ذهبي}}} & نص البيت من الشاطبية \\
\qw{كلمة مظللة} & الكلمة القرآنية المختلف فيها داخل البيت \\
\rz{Asim}{كلمة ملونة} & كلمة فيها رمز قارئ (لونها لون القارئ) \\
\colorbox{vocabback}{\textcolor{vocabframe}{\textbf{رمادي}}} & معاني الكلمات \\
\colorbox{white}{\textcolor{gold!80!black}{\textbf{ذهبي}}} & القراءات في البيت: من قرأ بماذا \\
\colorbox{ruleback}{\textcolor{ruleframe}{\textbf{أخضر}}} & قاعدة أو شرح اصطلاح \\
\colorbox{exback}{\textcolor{exframe}{\textbf{أزرق}}} & أمثلة من القرآن \\
\colorbox{tahrirback}{\textcolor{tahrirframe}{\textbf{بنفسجي}}} & تحريرات (من إتحاف البرية) \\
\colorbox{noteback}{\textcolor{noteframe}{\textbf{برتقالي}}} & تنبيه \\
\colorbox{sumback}{\textcolor{sumframe}{\textbf{خمري}}} & خلاصة الباب \\
\ayah{آية} & نص قرآني بالرسم العثماني، و\src{السورة: الآية} موضعه \\
\bottomrule
\end{tabular}
\end{center}

\begin{tanbih}[ملاحظة على التحريرات]
«التحريرات» هي ضبط ما يُقرأ به من طريق الشاطبية تحديدًا، ومنع خلط الطرق. وما نقلناه منها هو من منظومة «إتحاف البرية بتحريرات الشاطبية» للشيخ حسن خلف الحسيني رحمه الله. والأخذ بالتحريرات في الأداء يكون بالتلقي من شيخ متقن.
\end{tanbih}


\mainmatter
% ---------------- مقدمة الناظم (خطبة الكتاب) : الأبيات ١–٩٤ ----------------
% =====================================================================
%  مقدمة الناظم (١): الأبيات ١–١٩
%  البسملة والصلاة والحمد، ثم فضل القرآن وقارئه
% =====================================================================
\chapter{مقدمة الناظم (١): البسملة والحمد وفضل القرآن}

\begin{tanbih}[عن المقدمة]
مقدمة الشاطبية (وتسمى «خطبة الكتاب») ٩٤ بيتًا. ليس فيها خلاف بين القراء، بل هي: حمد وثناء، ثم فضل القرآن وأهله، ثم التعريف بالقراء السبعة ورواتهم، ثم \textbf{الرموز والاصطلاحات} (وهي أهم جزء لفهم ما بعده)، ثم وصايا وآداب.
لذلك نكتفي في أكثرها بمعاني الكلمات والمعنى المختصر، ونتوسع في أبيات الرموز (٤٤–٦٦).
\end{tanbih}

\section{البسملة والصلاة على النبي ﷺ والحمد}

\bayt{1}{بَدَأْتُ بِبِسْمِ اللهِ فِي النَّظْمِ أَوَّلَا}{تَبَارَكَ رَحْمَانًا رَحِيمًا وَمَوْئِلَا}
\begin{mufradat}
\kalima{النظم}{الشِّعر المنظوم، والمراد هنا هذه القصيدة.}
\kalima{تبارك}{كثُر خيره وعمّت بركته.}
\kalima{رحمانًا رحيمًا}{الرحمن: المنعم بجلائل النعم وعظائمها، والرحيم: المنعم بدقائقها.}
\kalima{مَوْئِلًا}{مرجعًا وملجأً يُلجأ إليه.}
\end{mufradat}
\maana بدأتُ قصيدتي بالبسملة أولًا، تبركًا باسم الله الذي كثرت بركته، فهو الرحمن الرحيم، وهو ملجأ اللاجئين.

\bayt{2}{وَثَنَّيْتُ صَلَّى اللهُ رَبِّي عَلَى الرِّضَا}{مُحَمَّدٍ الْمُهْدَى إِلَى النَّاسِ مُرْسَلَا}
\begin{mufradat}
\kalima{ثنّيت}{جعلتُه الثاني بعد البسملة.}
\kalima{الرِّضا}{المَرْضِيّ، أي الذي ارتضاه الله للنبوة.}
\kalima{المُهْدَى}{المُرسَل هديةً ورحمةً للناس.}
\end{mufradat}
\maana ثم ثنّيت بالصلاة على النبي محمد ﷺ الذي ارتضاه الله، وبعثه هدية لعباده، وواسطةً بينهم وبين خالقهم.

\bayt{3}{وَعِتْرَتِهِ ثُمَّ الصَّحَابَةِ ثُمَّ مَنْ}{تَلَاهُمْ عَلَى الْإِحْسَانِ بِالْخَيْرِ وُبَّلَا}
\begin{mufradat}
\kalima{عِتْرته}{أهل بيته وأقاربه الأدنَون.}
\kalima{تلاهم}{تبعهم.}
\kalima{وُبَّلًا}{جمع «وابل» وهو المطر الغزير.}
\end{mufradat}
\maana وصلى الله كذلك على آل النبي ﷺ ثم على صحابته ثم على من تبعهم بإحسان، صلاةً كثيرة تنزل عليهم بالخير كالمطر الغزير.

\bayt{4}{وَثَلَّثْتُ أَنَّ الْحَمْدَ لِلَّهِ دَائِمًا}{وَمَا لَيْسَ مَبْدُوءًا بِهِ أَجْذَمُ الْعَلَا}
\begin{mufradat}
\kalima{ثلّثت}{جعلتُه الثالث.}
\kalima{أجذم العلا}{مقطوع الرِّفعة والبركة. والأجذم: المقطوع.}
\end{mufradat}
\maana وجعلت الثالث: الحمد لله دائمًا؛ لأن كل أمر مهم لا يُبدأ بحمد الله فهو ناقص البركة. وفيه إشارة إلى الحديث: «كل أمر ذي بال لا يُبدأ فيه بالحمد لله فهو أجذم».

\section{فضل القرآن}

\bayt{5}{وَبَعْدُ فَحَبْلُ اللهِ فِينَا كِتَابُهُ}{فَجَاهِدْ بِهِ حِبْلَ الْعِدَا مُتَحَبِّلَا}
\begin{mufradat}
\kalima{حبل الله}{القرآن، سُمّي حبلًا لأنه سبب النجاة، والمتمسك به لا يضل.}
\kalima{حِبْل العِدا}{كيد الأعداء ومكرهم (الحِبْل بالكسر: الداهية).}
\kalima{متحبِّلًا}{متمسكًا به كمن يمسك بالحبل.}
\end{mufradat}
\maana أما بعد: فالقرآن هو حبل الله الممدود بيننا، فجاهد به كيد أعدائك متمسكًا به. وفيه إشارة إلى قوله تعالى: \ayah{وَٱعۡتَصِمُواْ بِحَبۡلِ ٱللَّهِ جَمِيعࣰا} \src{آل عمران: ١٠٣}.

\bayt{6}{وَأَخْلِقْ بِهِ إِذْ لَيْسَ يَخْلُقُ جِدَّةً}{جَدِيدًا مُوَالِيهِ عَلَى الْجِدِّ مُقْبِلَا}
\begin{mufradat}
\kalima{أَخْلِقْ به}{ما أجدره وما أحقّه (صيغة تعجب).}
\kalima{يَخْلُق}{يَبْلَى ويقدُم.}
\kalima{جِدّة}{الجِدَّة: كون الشيء جديدًا.}
\kalima{مواليه}{ملازمه ومحبه.}
\end{mufradat}
\maana ما أحقّ القرآن بأن يُتمسك به! فهو لا يبلى ولا تذهب جِدّته مهما كثر ترداده، وصاحبه الملازم له يبقى متجددًا مُقبِلًا على الجِدّ والاجتهاد.

\bayt{7}{وَقَارِئُهُ الْمَرْضِيُّ قَرَّ مِثَالُهُ}{كَالْأُتْرُجِّ حَالَيْهِ مُرِيحًا وَمُوكَلَا}
\begin{mufradat}
\kalima{قرّ مثاله}{ثبت مَثَلُه وصفته.}
\kalima{الأُتْرُجّ}{ثمرة طيبة الطعم والرائحة.}
\kalima{حاليه}{في حالتيه.}
\kalima{مُريحًا ومُوكَلًا}{ذا ريح طيبة، ومأكولًا طيب الطعم.}
\end{mufradat}
\maana قارئ القرآن العامل به الذي يرضاه الله ثبت مَثَلُه في الحديث بأنه كالأترجة: طيبة في حالتيها، ريحها طيب وطعمها طيب. وهذا من قوله ﷺ: «مثل المؤمن الذي يقرأ القرآن كمثل الأترجة: ريحها طيب وطعمها طيب» (متفق عليه).

\bayt{8}{هُوَ الْمُرْتَضَى أَمًّا إِذَا كَانَ أُمَّةً}{وَيَمَّمَهُ ظِلُّ الرَّزَانَةِ قَنْقَلَا}
\begin{mufradat}
\kalima{أَمًّا}{مقصودًا يُقتدى به (من أَمَّ يؤُمّ: قصد).}
\kalima{أُمّة}{رجلًا جامعًا لخصال الخير، كقوله تعالى: \ayah{إِنَّ إِبۡرَٰهِيمَ كَانَ أُمَّةࣰ} \src{النحل: ١٢٠}.}
\kalima{يمّمه}{قصده.}
\kalima{الرزانة}{الوقار والثبات.}
\kalima{قَنْقَلًا}{كثيرًا وافرًا (وأصل القنقل: المكيال الضخم).}
\end{mufradat}
\maana هذا القارئ هو المرضي الذي يُقصد ويُقتدى به، إذا جمع خصال الخير وأقبل عليه الوقار والسكينة بكثرة.

\bayt{9}{هُوَ الْحُرُّ إِنْ كَانَ الْحَرِيَّ حَوَارِيًا}{لَهُ بِتَحَرِّيهِ إِلَى أَنْ تَنَبَّلَا}
\begin{mufradat}
\kalima{الحُرّ}{الكريم الذي لم يستعبده هواه.}
\kalima{الحَرِيّ}{الجدير والحقيق.}
\kalima{حواريًّا}{ناصرًا مخلصًا (كحواريّي الأنبياء).}
\kalima{تحرّيه}{طلبه الأحرى والأصوب.}
\kalima{تنبّلا}{صار نبيلًا فاضلًا.}
\end{mufradat}
\maana وهو الكريم حقًّا إن كان جديرًا بأن يكون ناصرًا مخلصًا للقرآن، يتحرّى العمل به ويقصد الصواب فيه حتى يبلغ النبل والفضل.

\bayt{10}{وَإِنَّ كِتَابَ اللهِ أَوْثَقُ شَافِعٍ}{وَأَغْنَى غَنَاءٍ وَاهِبًا مُتَفَضِّلَا}
\begin{mufradat}
\kalima{أوثق شافع}{أوثق من يشفع، أي شفاعته مقبولة.}
\kalima{أغنى غَناء}{أعظم نفعًا وكفاية.}
\kalima{واهبًا متفضلًا}{معطيًا الأجر بلا مقابل.}
\end{mufradat}
\maana القرآن أوثق شفيع لصاحبه يوم القيامة، وأعظم ما ينفعه ويغنيه، يهب له الثواب تفضلًا. وفيه إشارة إلى حديث: «اقرؤوا القرآن فإنه يأتي يوم القيامة شفيعًا لأصحابه» (رواه مسلم).

\bayt{11}{وَخَيْرُ جَلِيسٍ لَا يُمَلُّ حَدِيثُهُ}{وَتَرْدَادُهُ يَزْدَادُ فِيهِ تَجَمُّلَا}
\begin{mufradat}
\kalima{ترداده}{تكراره.}
\kalima{تجمّلًا}{حُسنًا وجمالًا.}
\end{mufradat}
\maana والقرآن خير جليس، لا يُملّ حديثه، وكلما كرّرته ازداد حسنًا وجمالًا في نفسك، بخلاف كلام البشر الذي يُملّ مع التكرار.

\bayt{12}{وَحَيْثُ الْفَتَى يَرْتَاعُ فِي ظُلُمَاتِهِ}{مِنَ الْقَبْرِ يَلْقَاهُ سَنًا مُتَهَلِّلَا}
\begin{mufradat}
\kalima{يرتاع}{يفزع ويخاف.}
\kalima{سَنًا}{ضوءًا ونورًا.}
\kalima{متهللًا}{مشرقًا مستبشرًا.}
\end{mufradat}
\maana وحين يفزع الإنسان في ظلمات القبر، يلقاه القرآن نورًا مشرقًا يؤنسه.

\bayt{13}{هُنَالِكَ يَهْنِيهِ مَقِيلًا وَرَوْضَةً}{وَمِنْ أَجْلِهِ فِي ذِرْوَةِ الْعِزِّ يَجْتَلِي}
\begin{mufradat}
\kalima{يَهْنيه}{يُهنِّئه (أصلها يهنئه بالهمز فخُففت).}
\kalima{مَقِيلًا}{مكان الراحة وقت القيلولة، والمراد القبر.}
\kalima{ذِرْوة العز}{أعلى درجات العزة.}
\kalima{يجتلي}{يظهر ويُرى في أبهى صورة.}
\end{mufradat}
\maana هناك في القبر يهنئه القرآن بمكان راحة وروضة من رياض الجنة، وبسببه يظهر صاحبه في أعلى منازل العزة.

\bayt{14}{يُنَاشِدُ فِي إِرْضَائِهِ لِحَبِيبِهِ}{وَأَجْدِرْ بِهِ سُؤْلًا إِلَيْهِ مُوَصَّلَا}
\begin{mufradat}
\kalima{يناشد}{يسأل ربه ويلحّ في السؤال.}
\kalima{لحبيبه}{أي لقارئه الذي أحبه.}
\kalima{أجدِرْ به}{ما أحقّه (تعجب).}
\kalima{سؤلًا موصّلًا}{طلبًا مقبولًا مُبلَّغًا.}
\end{mufradat}
\maana يسأل القرآن ربه أن يرضى عن قارئه، وما أحقّ هذا الطلب بالقبول. وفيه إشارة إلى حديث: «يجيء القرآن يوم القيامة فيقول: يا رب حَلِّه، فيُلبس تاج الكرامة... ثم يقول: يا رب ارضَ عنه، فيرضى عنه» (رواه الترمذي).

\bayt{15}{فَيَا أَيُّهَا الْقَارِي بِهِ مُتَمَسِّكًا}{مُجِلًّا لَهُ فِي كُلِّ حَالٍ مُبَجِّلَا}
\begin{mufradat}
\kalima{مُجِلًّا مُبجِّلًا}{معظّمًا مكرّمًا.}
\end{mufradat}
\maana يا أيها القارئ المتمسك بالقرآن العامل به، المعظم له في كل أحوالك (وجواب النداء في البيت التالي).

\bayt{16}{هَنِيئًا مَرِيئًا وَالِدَاكَ عَلَيْهِمَا}{مَلَابِسُ أَنْوَارٍ مِنَ التَّاجِ وَالْحُلَى}
\begin{mufradat}
\kalima{هنيئًا مريئًا}{دعاء بأن يكون الخير سائغًا لا تعب فيه.}
\kalima{الحُلَى}{جمع حِلية، وهي الزينة.}
\end{mufradat}
\maana هنيئًا لك ما أعدّه الله لك، حتى إن والديك يُلبَسان يوم القيامة تاجًا وحُللًا من نور. وفيه إشارة إلى حديث: «من قرأ القرآن وعمل بما فيه أُلبس والداه تاجًا يوم القيامة ضوؤه أحسن من ضوء الشمس...» (رواه أبو داود).

\bayt{17}{فَمَا ظَنُّكُمْ بِالنَّجْلِ عِنْدَ جَزَائِهِ}{أُولَئِكَ أَهْلُ اللهِ وَالصَّفْوَةُ الْمَلَا}
\begin{mufradat}
\kalima{النَّجْل}{الولد، والمراد القارئ نفسه.}
\kalima{الصَّفْوة}{الخيار.}
\kalima{المَلَا}{الأشراف.}
\end{mufradat}
\maana إذا كان هذا إكرام الوالدين، فكيف يكون إكرام الولد القارئ نفسه؟! إنهم أهل الله وخاصته، كما في الحديث: «أهل القرآن هم أهل الله وخاصته» (رواه ابن ماجه).

\bayt{18}{أُولُو الْبِرِّ وَالْإِحْسَانِ وَالصَّبْرِ وَالتُّقَى}{حُلَاهُمْ بِهَا جَاءَ الْقُرَانُ مُفَصَّلَا}
\begin{mufradat}
\kalima{حُلاهم}{صفاتهم الحسنة.}
\end{mufradat}
\maana وهم أصحاب البر والإحسان والصبر والتقوى، وهذه صفاتهم قد جاء القرآن بتفصيلها في آيات كثيرة.

\bayt{19}{عَلَيْكَ بِهَا مَا عِشْتَ فِيهَا مُنَافِسًا}{وَبِعْ نَفْسَكَ الدُّنْيَا بِأَنْفَاسِهَا الْعُلَى}
\begin{mufradat}
\kalima{عليك بها}{الزَمْ هذه الصفات.}
\kalima{منافسًا}{مسابقًا غيرك راغبًا فيها.}
\kalima{أنفاسها العُلى}{النفيس الغالي من الصفات العالية.}
\end{mufradat}
\maana الزم هذه الصفات طول حياتك متسابقًا فيها، واجعل الدنيا ثمنًا تبذله لتنال نفسُك تلك المراتب العالية، أي قدّم الآخرة على الدنيا.
       % ١–١٩   : الحمد وفضل القرآن وأهله
% =====================================================================
%  مقدمة الناظم (٢): الأبيات ٢٠–٤٣ — القراء السبعة ورواتهم
% =====================================================================
\chapter{مقدمة الناظم (٢): القراء السبعة ورواتهم}

\begin{tanbih}[فكرة هذا القسم]
يعرّفك الناظم هنا بالقراء السبعة، ولكل قارئ راويان. حفظ هذه الأسماء بترتيبها \textbf{مهم جدًّا}؛ لأن رموز «أبجد» مبنية على هذا الترتيب نفسه.
ولوّنا اسم كل قارئ وراوييه بلونه في الأبيات.
\end{tanbih}

\bayt{20}{جَزَى اللهُ بِالْخَيْرَاتِ عَنَّا أَئِمَّةً}{لَنَا نَقَلُوا الْقُرْآنَ عَذْبًا وَسَلْسَلَا}
\begin{mufradat}
\kalima{عذبًا وسلسلًا}{حلوًا سهلًا على اللسان.}
\end{mufradat}
\maana دعاء للأئمة الذين نقلوا إلينا القرآن نقلًا صحيحًا متقنًا، سهلًا عذبًا كما أُنزل.

\bayt{21}{فَمِنْهُمْ بُدُورٌ سَبْعَةٌ قَدْ تَوَسَّطَتْ}{سَمَاءَ الْعُلَى وَالْعَدْلِ زُهْرًا وَكُمَّلَا}
\begin{mufradat}
\kalima{بدور}{جمع بدر، والمراد القراء السبعة.}
\kalima{زُهرًا}{مضيئة مشرقة.}
\kalima{كُمَّلًا}{جمع كامل.}
\end{mufradat}
\maana من هؤلاء الأئمة سبعة كالبدور، توسطوا سماء الرفعة والعدالة، مشرقين كاملين.

\bayt{22}{لَهَا شُهُبٌ عَنْهَا اسْتَنَارَتْ فَنَوَّرَتْ}{سَوَادَ الدُّجَى حَتَّى تَفَرَّقَ وَانْجَلَى}
\begin{mufradat}
\kalima{شُهُب}{جمع شهاب، وهو النجم المضيء، والمراد الرواة.}
\kalima{الدُّجى}{الظلام.}
\end{mufradat}
\maana ولهؤلاء البدور (القراء) نجوم (رواة) أخذوا النور منهم، فأضاؤوا ظلام الجهل حتى انكشف.

\bayt{23}{وَسَوْفَ تَرَاهُمْ وَاحِدًا بَعْدَ وَاحِدٍ}{مَعَ اثْنَيْنِ مِنْ أَصْحَابِهِ مُتَمَثِّلَا}
\begin{mufradat}
\kalima{متمثلًا}{ظاهرًا ماثلًا أمامك.}
\end{mufradat}
\maana سأذكر لك القراء واحدًا واحدًا، ومع كل قارئ راويين من أصحابه.

\bayt{24}{تَخَيَّرَهُمْ نُقَّادُهُمْ كُلَّ بَارِعٍ}{وَلَيْسَ عَلَى قُرْآنِهِ مُتَأَكِّلَا}
\begin{mufradat}
\kalima{نُقّادهم}{العلماء الذين ميّزوا الرواة وانتقوهم.}
\kalima{بارع}{فائق في علمه.}
\kalima{متأكّلًا}{يجعل القرآن وسيلة للتكسّب.}
\end{mufradat}
\maana اختار العلماءُ النقادُ هؤلاء الرواة؛ فكل واحد منهم بارع متقن، ولا يتخذ القرآن وسيلة للأكل والكسب.

\section{نافع المدني وراوياه}
\bayt{25}{فَأَمَّا الْكَرِيمُ السِّرِّ فِي الطِّيبِ \qname{Nafi}{نَافِعٌ}}{فَذَاكَ الَّذِي اخْتَارَ الْمَدِينَةَ مَنْزِلَا}
\begin{mufradat}
\kalima{الكريم السر}{طاهر الباطن.}
\kalima{في الطِّيب}{يُروى أنه كان إذا تكلم شُمّت من فمه رائحة المسك، فسئل عن ذلك فقال: رأيت النبي ﷺ في المنام يقرأ في فيّ.}
\end{mufradat}
\maana القارئ الأول: \Nafi{} بن عبد الرحمن بن أبي نُعيم، اختار المدينة المنورة دارًا ومقامًا (ت ١٦٩ هـ).

\bayt{26}{\qname{Nafi}{وَقَالُونُ} عِيسَى ثُمَّ عُثْمَانُ \qname{Nafi}{وَرْشُهُمْ}}{بِصُحْبَتِهِ الْمَجْدَ الرَّفِيعَ تَأَثَّلَا}
\begin{mufradat}
\kalima{تأثّلا}{اكتسبا وأصّلا.}
\end{mufradat}
\maana راويا \Nafi: \Qalun{} واسمه \textbf{عيسى} (ت ٢٢٠ هـ)، ثم \Warsh{} واسمه \textbf{عثمان} (ت ١٩٧ هـ)، وقد نالا المجد الرفيع بصحبتهما \Nafi.

\section{ابن كثير المكي وراوياه}
\bayt{27}{وَمَكَّةُ عَبْدُ اللهِ فِيهَا مُقَامُهُ}{هُوَ \qname{IbnKathir}{ابْنُ كَثِيرٍ} كَاثِرُ الْقَوْمِ مُعْتَلَى}
\begin{mufradat}
\kalima{كاثر القوم}{غالبهم وفائقهم (وفيه تلاعب لطيف باسم «كثير»).}
\kalima{مُعتلى}{عاليًا رفيعًا.}
\end{mufradat}
\maana القارئ الثاني: \textbf{عبد الله} \IbnKathir{} المكي، أقام بمكة، وفاق قومه علمًا ورفعة (ت ١٢٠ هـ).

\bayt{28}{رَوَى أَحْمَدُ \qname{IbnKathir}{الْبَزِّي} لَهُ وَمُحَمَّدٌ}{عَلَى سَنَدٍ وَهْوَ الْمُلَقَّبُ \qname{IbnKathir}{قُنْبُلَا}}
\begin{mufradat}
\kalima{على سند}{أي رويا عنه بواسطة رجال، لا مباشرة.}
\end{mufradat}
\maana راويا \IbnKathir: \Bazzi{} واسمه \textbf{أحمد} (ت ٢٥٠ هـ)، و\Qunbul{} واسمه \textbf{محمد} (ت ٢٩١ هـ)، وكلاهما روى عنه بواسطة.

\section{أبو عمرو البصري وراوياه}
\bayt{29}{وَأَمَّا الْإِمَامُ الْمَازِنِيُّ صَرِيحُهُمْ}{\qname{AbuAmr}{أَبُو عَمْرٍو الْبَصْرِي} فَوَالِدُهُ الْعَلَا}
\begin{mufradat}
\kalima{المازني}{نسبة إلى بني مازن من تميم.}
\kalima{صريحهم}{خالص النسب العربي فيهم.}
\end{mufradat}
\maana القارئ الثالث: \AbuAmr{} البصري، عربي صريح النسب من بني مازن، وأبوه اسمه \textbf{العلاء} (ت ١٥٤ هـ). ولذلك سيسميه الناظم كثيرًا «ابن العلاء» و«فتى العلا».

\bayt{30}{أَفَاضَ عَلَى يَحْيَى الْيَزِيدِيِّ سَيْبَهُ}{فَأَصْبَحَ بِالْعَذْبِ الْفُرَاتِ مُعَلَّلَا}
\begin{mufradat}
\kalima{سَيْبَه}{عطاءه، أي علمه.}
\kalima{العذب الفرات}{الماء الحلو الشديد العذوبة.}
\kalima{معلَّلًا}{مسقيًّا مرة بعد مرة.}
\end{mufradat}
\maana أفاض \AbuAmr{} علمه على تلميذه \textbf{يحيى اليزيدي}، فارتوى منه، وعن اليزيدي أخذ الراويان المذكوران بعده.

\bayt{31}{أَبُو عُمَرَ \qname{AbuAmr}{الدُّورِي} وَصَالِحُهُمْ أَبُو}{شُعَيْبٍ هُوَ \qname{AbuAmr}{السُّوسِيُّ} عَنْهُ تَقَبَّلَا}
\maana راويا \AbuAmr{} (بواسطة اليزيدي): \qname{AbuAmr}{الدوري} حفص بن عمر، كنيته \textbf{أبو عُمَر} (ت ٢٤٦ هـ)، و\Susi{} صالح بن زياد، كنيته \textbf{أبو شعيب} (ت ٢٦١ هـ).

\section{ابن عامر الشامي وراوياه}
\bayt{32}{وَأَمَّا دِمَشْقُ الشَّامِ دَارُ \qname{IbnAmir}{ابْنِ عَامِرٍ}}{فَتِلْكَ بِعَبْدِ اللهِ طَابَتْ مُحَلَّلَا}
\begin{mufradat}
\kalima{مُحَلَّلًا}{منزلًا ومكانَ حلول.}
\end{mufradat}
\maana القارئ الرابع: \textbf{عبد الله} \IbnAmir{} اليَحْصُبي، داره دمشق، وطابت به منزلًا (ت ١١٨ هـ). وهو أعلى القراء السبعة سندًا.

\bayt{33}{\qname{IbnAmir}{هِشَامٌ} وَعَبْدُ اللهِ وَهْوَ انْتِسَابُهُ}{\qname{IbnAmir}{لِذَكْوَانَ} بِالْإِسْنَادِ عَنْهُ تَنَقَّلَا}
\maana راويا \IbnAmir: \Hisham{} بن عمار (ت ٢٤٥ هـ)، و\textbf{عبد الله} بن أحمد \IbnDhakwan{} — نُسب إلى جده ذكوان (ت ٢٤٢ هـ). وكلاهما روى عنه بواسطة.

\section{القراء الكوفيون الثلاثة}
\bayt{34}{وَبِالْكُوفَةِ الْغَرَّاءِ مِنْهُمْ ثَلَاثَةٌ}{أَذَاعُوا فَقَدْ ضَاعَتْ شَذًا وَقَرَنْفُلَا}
\begin{mufradat}
\kalima{الغرّاء}{البيضاء المشهورة.}
\kalima{أذاعوا}{نشروا القراءة.}
\kalima{ضاعت}{فاحت رائحتها.}
\kalima{شذًا}{الرائحة الطيبة.}
\end{mufradat}
\maana وفي الكوفة ثلاثة من السبعة نشروا القرآن، ففاحت الكوفة بعلمهم كرائحة الطيب والقرنفل. وهؤلاء الثلاثة هم «الكوفيون» ورمزهم \rzlg{ث}.

\bayt{35}{فَأَمَّا أَبُو بَكْرٍ وَ\qname{Asim}{عَاصِمٌ} اسْمُهُ}{\qname{Asim}{فَشُعْبَةُ} رَاوِيهِ الْمُبَرِّزُ أَفْضَلَا}
\begin{mufradat}
\kalima{المبرِّز}{المتقدم الفائق.}
\end{mufradat}
\maana القارئ الخامس: \Asim{} بن أبي النَّجود، كنيته \textbf{أبو بكر} (ت ١٢٧ هـ). وراويه الأول \Shuba، المتقدم الفاضل.

\bayt{36}{وَذَاكَ ابْنُ عَيَّاشٍ أَبُو بَكْرٍ الرِّضَا}{وَ\qname{Asim}{حَفْصٌ} وَبِالْإِتْقَانِ كَانَ مُفَضَّلَا}
\maana \Shuba{} هو ابن عيّاش، وكنيته أيضًا \textbf{أبو بكر} (ت ١٩٣ هـ). والراوي الثاني \Hafs{} بن سليمان، وكان مفضّلًا بالإتقان (ت ١٨٠ هـ). وروايته هي المنتشرة اليوم في أكثر العالم الإسلامي.

\bayt{37}{وَ\qname{Hamza}{حَمْزَةُ} مَا أَزْكَاهُ مِنْ مُتَوَرِّعٍ}{إِمَامًا صَبُورًا لِلْقُرَانِ مُرَتِّلَا}
\maana القارئ السادس: \Hamza{} بن حبيب الزيّات، ما أطهره من ورِع! كان إمامًا صبورًا مرتّلًا للقرآن (ت ١٥٦ هـ).

\bayt{38}{رَوَى \qname{Hamza}{خَلَفٌ} عَنْهُ وَ\qname{Hamza}{خَلَّادٌ} الَّذِي}{رَوَاهُ سُلَيْمٌ مُتْقَنًا وَمُحَصَّلَا}
\maana راويا \Hamza: \Khalaf{} بن هشام (ت ٢٢٩ هـ) و\Khallad{} بن خالد (ت ٢٢٠ هـ)، وكلاهما أخذ عن \textbf{سُلَيم} بن عيسى عن \Hamza.

\bayt{39}{وَأَمَّا عَلِيٌّ \qname{Kisai}{فَالْكِسَائِيُّ} نَعْتُهُ}{لِمَا كَانَ فِي الْإِحْرَامِ فِيهِ تَسَرْبَلَا}
\begin{mufradat}
\kalima{تسربل}{لبس.}
\end{mufradat}
\maana القارئ السابع: \textbf{علي} بن حمزة، ولقبه \Kisai؛ لأنه أحرم في كساء (ت ١٨٩ هـ).

\bayt{40}{رَوَى لَيْثُهُمْ عَنْهُ \qname{Kisai}{أَبُو الْحَارِثِ} الرِّضَا}{وَحَفْصٌ هُوَ \qname{Kisai}{الدُّورِي} وَفِي الذِّكْرِ قَدْ خَلَا}
\begin{mufradat}
\kalima{ليثهم}{اسمه الليث بن خالد.}
\kalima{في الذكر قد خلا}{سبق ذكره (مع أبي عمرو).}
\end{mufradat}
\maana راويا \Kisai: \AbulHarith{} الليث بن خالد (ت ٢٤٠ هـ)، و\textbf{حفص الدوري}، وهو نفسه راوي \AbiAmr{} الذي سبق ذكره.

\begin{tanbih}[انتبه: الدوري راوٍ لقارئين]
الدوري (حفص بن عمر) روى عن \AbiAmr{} وروى أيضًا عن \Kisai. فإذا قيل «الدوري» في باب فانظر: هل الكلام عن \AbiAmr{} أم \Kisai؟ ولهذا جعلنا لكل منهما لون قارئه.
\end{tanbih}

\section{أنساب القراء وطرق الرواية}
\bayt{41}{\qname{AbuAmr}{أَبُو عَمْرِهِمْ} وَالْيَحْصَبِيُّ \qname{IbnAmir}{ابْنُ عَامِرٍ}}{صَرِيحٌ وَبَاقِيهِمْ أَحَاطَ بِهِ الْوَلَا}
\begin{mufradat}
\kalima{صريح}{عربي خالص النسب.}
\kalima{الوَلا}{الولاء، أي أنهم من الموالي (أصولهم غير عربية).}
\end{mufradat}
\maana \AbuAmr{} و\IbnAmir{} عربيان صريحا النسب، والباقون من الموالي. (وهذا لا ينقص من قدرهم شيئًا، فالفضل بالعلم والتقوى.)

\bayt{42}{لَهُمْ طُرُقٌ يُهْدَى بِهَا كُلُّ طَارِقٍ}{وَلَا طَارِقٌ يُخْشَى بِهَا مُتَمَحِّلَا}
\begin{mufradat}
\kalima{طُرُق}{طرق الرواية والأسانيد.}
\kalima{طارق}{قاصد (من يأتي ليلًا).}
\kalima{ولا طارقٌ يُخشى}{ولا داهية أو ضلال يُخاف منه.}
\kalima{متمحِّلًا}{ماكرًا كائدًا.}
\end{mufradat}
\maana لهؤلاء الرواة طرق (أسانيد) يهتدي بها كل طالب، ولا يُخشى فيها ضلال.

\begin{qaida}[ما الفرق بين القراءة والرواية والطريق؟]
\textbf{القراءة}: ما يُنسب إلى الإمام (مثل: قراءة \Nafi).\\
\textbf{الرواية}: ما يُنسب إلى الراوي عن الإمام (مثل: رواية \Qalun{} عن \Nafi).\\
\textbf{الطريق}: ما يُنسب إلى من أخذ عن الراوي وإن نزل (مثل: طريق أبي نشيط عن \Qalun، وهو طريق الشاطبية).\\
والشاطبية نَظْمٌ لكتاب «التيسير» للداني، فكل راوٍ فيها من طريق واحد غالبًا.
\end{qaida}

\bayt{43}{وَهُنَّ اللَّوَاتِي لِلْمُوَاتِي نَصَبْتُهَا}{مَنَاصِبَ فَانْصَبْ فِي نِصَابِكَ مُفْضِلَا}
\begin{mufradat}
\kalima{المُواتي}{الموافق المطاوع، أي الطالب المجتهد.}
\kalima{مناصب}{علامات منصوبة يُهتدى بها.}
\kalima{فانْصَبْ}{اجتهد واتعب.}
\kalima{نِصابك}{أصلك وحظك.}
\kalima{مُفضِلًا}{محسنًا.}
\end{mufradat}
\maana هذه الطرق نصبتُها علامات للطالب الموافق، فاجتهد في طلب هذا العلم محسنًا.
      % ٢٠–٤٣  : القراء السبعة ورواتهم
% =====================================================================
%  مقدمة الناظم (٣): الأبيات ٤٤–٦٦ — الرموز والاصطلاحات
%  (أهم جزء في المقدمة: عليه يُبنى فهم بقية القصيدة)
% =====================================================================
\chapter{مقدمة الناظم (٣): الرموز والاصطلاحات}

\begin{tanbih}[لماذا هذا القسم مهم؟]
الشاطبي لم يذكر أسماء القراء صراحة في أكثر القصيدة، بل وضع لكل قارئ \textbf{حرفًا} يرمز إليه، ولبعض الجماعات \textbf{حرفًا أو كلمة}، ووضع قواعد يُفهم منها قراءة من لم يُذكر. فمن لم يفهم هذه الأبيات لم يفهم القصيدة.
\end{tanbih}

\section{الرموز الحرفية لكل قارئ وراوٍ}

\bayt{44}{وَهَا أَنَا ذَا أَسْعَى لَعَلَّ حُرُوفَهُمْ}{يَطُوعُ بِهَا نَظْمُ الْقَوَافِي مُسَهَّلَا}
\begin{mufradat}
\kalima{حروفهم}{قراءاتهم (الحرف: القراءة).}
\kalima{يطوع}{ينقاد.}
\end{mufradat}
\maana أسعى راجيًا أن تنقاد قراءاتهم لنظم القصيدة بسهولة.

\bayt{45}{جَعَلْتُ أَبَا جَادٍ عَلَى كُلِّ قَارِئٍ}{دَلِيلًا عَلَى الْمَنْظُومِ أَوَّلَ أَوَّلَا}
\begin{mufradat}
\kalima{أبا جاد}{حروف «أبجد هوز حطي كلمن سعفص قرشت».}
\kalima{أوّلَ أوّلا}{بالترتيب: الأول للأول، والذي بعده لمن بعده.}
\end{mufradat}
\maana جعلت حروف «أبجد» رموزًا للقراء والرواة بالترتيب الذي ذكرتهم به: الحرف الأول للقارئ الأول، والثاني لراويه الأول، والثالث لراويه الثاني، وهكذا.

\begin{qaida}[جدول الرموز الحرفية]
\centering
\renewcommand{\arraystretch}{1.5}
\begin{tabular}{c c c}
\toprule
\textbf{الإمام} & \textbf{الراوي الأول} & \textbf{الراوي الثاني}\\
\midrule
\rzl{Nafi}{أ} \Nafi & \rzl{Nafi}{ب} \Qalun & \rzl{Nafi}{ج} \Warsh\\
\rzl{IbnKathir}{د} \IbnKathir & \rzl{IbnKathir}{هـ} \Bazzi & \rzl{IbnKathir}{ز} \Qunbul\\
\rzl{AbuAmr}{ح} \AbuAmr & \rzl{AbuAmr}{ط} \qname{AbuAmr}{الدوري} & \rzl{AbuAmr}{ي} \Susi\\
\rzl{IbnAmir}{ك} \IbnAmir & \rzl{IbnAmir}{ل} \Hisham & \rzl{IbnAmir}{م} \IbnDhakwan\\
\rzl{Asim}{ن} \Asim & \rzl{Asim}{ص} \Shuba & \rzl{Asim}{ع} \Hafs\\
\rzl{Hamza}{ف} \Hamza & \rzl{Hamza}{ض} \Khalaf & \rzl{Hamza}{ق} \Khallad\\
\rzl{Kisai}{ر} \Kisai & \rzl{Kisai}{س} \AbulHarith & \rzl{Kisai}{ت} \qname{Kisai}{الدوري}\\
\bottomrule
\end{tabular}

\smallskip
\raggedleft
\textbf{تذكّر}: «أبج — دهز — حطي — كلم — نصع — فضق — رست».
\end{qaida}

\bayt{46}{وَمِنْ بَعْدِ ذِكْرِي الْحَرْفَ أُسْمِي رِجَالَهُ}{مَتَى تَنْقَضِي آتِيكَ بِالْوَاوِ فَيْصَلَا}
\begin{mufradat}
\kalima{الحرف}{الكلمة القرآنية المختلف فيها.}
\kalima{أُسْمي رجاله}{أذكر رموز قرائه.}
\kalima{فيصلًا}{فاصلًا.}
\end{mufradat}
\maana عندي قاعدتان: (١) أذكر الكلمة القرآنية \textbf{أولًا} ثم أذكر رمز من قرأها \textbf{بعدها}. (٢) إذا انتهت المسألة أتيت بـ\textbf{واو} تفصلها عن المسألة التالية.

\begin{amthila}[مثال تطبيقي من البيت ١٠٨]
\begin{center}
\qw{وَمَالِكِ} يَوْمِ الدِّينِ \rz{Kisai}{رَاوِيهِ} \rz{Asim}{نَاصِرٌ} \quad \textbf{وَ}عِنْدَ سِرَاطٍ ...
\end{center}
\begin{itemize}
\item الكلمة القرآنية: «مالك» \ila{} ثم الرمز بعدها: \rz{Kisai}{ر} (\Kisai) و\rz{Asim}{ن} (\Asim).
\item ثم جاءت \textbf{الواو} في «وعند» فعرفنا أن المسألة انتهت وبدأت مسألة جديدة.
\end{itemize}
\end{amthila}

\bayt{47}{سِوَى أَحْرُفٍ لَا رِيبَةٌ فِي اتِّصَالِهَا}{وَبِاللَّفْظِ أَسْتَغْنِي عَنِ الْقَيْدِ إِنْ جَلَا}
\begin{mufradat}
\kalima{لا ريبة}{لا شك.}
\kalima{القيد}{الوصف المبيِّن للقراءة (كقوله: بالمد، بالفتح...).}
\kalima{جلا}{ظهر واتضح.}
\end{mufradat}
\maana
\textbf{الأول}: أنه قد يترك واو الفصل في مواضع قليلة لا يشك أحد أن الكلام فيها متصل بما قبله أو منفصل عنه.
\textbf{الثاني}: إذا كان نطقه بالكلمة يوضح القراءة فإنه لا يقيّدها بوصف.

\begin{amthila}
\begin{itemize}
\item \textbf{الاستغناء باللفظ}: قال «ومالك» بالألف، فعرفنا أن \Kisai{} و\Asim{} يقرآن بالألف دون أن يقول «بالمد».
\item \textbf{حرف لا ريبة فيه}: قال في البيت ١١٥: «... بالضم شمللا \textbf{كَمَا} بهم الأسباب»، فالكاف في «كما» ليست رمزًا لابن عامر، لأنها واضحة أنها كاف التشبيه قبل مثال.
\end{itemize}
\end{amthila}

\bayt{48}{وَرُبَّ مَكَانٍ كَرَّرَ الْحَرْفَ قَبْلَهَا}{لِمَا عَارِضٍ وَالْأَمْرُ لَيْسَ مُهَوِّلَا}
\begin{mufradat}
\kalima{لما عارضٍ}{لسبب عارض، كضرورة الوزن.}
\kalima{ليس مهوِّلًا}{ليس مخيفًا ولا صعبًا.}
\end{mufradat}
\maana قد يكرر الرمز (يأتي بالرمز مرتين للقارئ نفسه) قبل واو الفصل لضرورة الوزن أو لغيرها، وهذا لا يضر ولا يُشكل؛ لأن المعنى واحد.

\section{رموز الجماعة الحرفية: ث خ ذ ظ غ ش}

\begin{tanbih}
بعد «أبجد ... قرشت» بقيت ستة أحرف: «ثَخِذْ ظَغِشْ». جعل كل حرف منها رمزًا لجماعة من القراء.
\end{tanbih}

\bayt{49}{وَمِنْهُنَّ لِلْكُوفِي \rzg{ثَاءٌ} مُثَلَّثٌ}{وَسِتَّتُهُمْ \rzg{بِالْخَاءِ} لَيْسَ بِأَغْفَلَا}
\begin{mufradat}
\kalima{ثاء مثلث}{الثاء ذات النقط الثلاث.}
\kalima{ليس بأغفلا}{ليس خاليًا من النقط (الأغفل: الحرف غير المنقوط)، أي الخاء المعجمة لا الحاء.}
\end{mufradat}

\bayt{50}{عَنَيْتُ الْأُولَى أَثْبَتُّهُمْ بَعْدَ نَافِعٍ}{وَكُوفٍ وَشَامٍ \rzg{ذَالُهُمْ} لَيْسَ مُغْفَلَا}
\begin{mufradat}
\kalima{عنيت}{قصدت.}
\kalima{الأُولى أثبتّهم بعد نافع}{أي القراء الستة الذين ذكرتهم بعد نافع.}
\end{mufradat}

\bayt{51}{وَكُوفٍ مَعَ الْمَكِّيِّ \rzg{بِالظَّاءِ} مُعْجَمًا}{وَكُوفٍ وَبَصْرٍ \rzg{غَيْنُهُمْ} لَيْسَ مُهْمَلَا}
\begin{mufradat}
\kalima{معجمًا / ليس مهملًا}{منقوطًا. (المعجم: المنقوط، والمهمل: غير المنقوط.)}
\end{mufradat}

\bayt{52}{وَذُو النَّقْطِ \rzg{شِينٌ} لِلْكِسَائِي وَحَمْزَةٍ}{وَقُلْ فِيهِمَا مَعْ شُعْبَةٍ \rzg{صُحْبَةٌ} تَلَا}

\maana[(الأبيات ٤٩–٥٢)]
\begin{itemize}
\item \rzlg{ث} = الكوفيون الثلاثة: \Asim{} و\Hamza{} و\Kisai.
\item \rzlg{خ} = القراء الستة غير \Nafi: \IbnKathir، \AbuAmr، \IbnAmir، \Asim، \Hamza، \Kisai.
\item \rzlg{ذ} = الكوفيون و\IbnAmir{} الشامي.
\item \rzlg{ظ} = الكوفيون و\IbnKathir{} المكي.
\item \rzlg{غ} = الكوفيون و\AbuAmr{} البصري.
\item \rzlg{ش} = \Hamza{} و\Kisai.
\end{itemize}
ثم بدأ بالرموز الكلمية: فقال إن كلمة \rzg{صُحْبة} رمز لـ\Hamza{} و\Kisai{} ومعهما \Shuba.
وأكّد في الأبيات أن هذه الحروف \textbf{منقوطة} (الخاء لا الحاء، والذال لا الدال، والظاء لا الطاء، والغين لا العين، والشين لا السين) حتى لا تلتبس بحروف «أبجد».

\section{رموز الجماعة الكلمية}

\bayt{53}{\rzg{صِحَابٌ} هُمَا مَعْ حَفْصِهِمْ \rzg{عَمَّ} نَافِعٌ}{وَشَامٍ \rzg{سَمَا} فِي نَافِعٍ وَفَتَى الْعَلَا}
\bayt{54}{وَمَكٍّ وَ\rzg{حَقٌّ} فِيهِ وَابْنِ الْعَلَاءِ قُلْ}{وَقُلْ فِيهِمَا وَالْيَحْصَبِي \rzg{نَفَرٌ} حَلَا}
\bayt{55}{\rzg{وَحِرْمِيٌّ} الْمَكِّيُّ فِيهِ وَنَافِعٌ}{\rzg{وَحِصْنٌ} عَنِ الْكُوفِي وَنَافِعِهِمْ عَلَا}
\begin{mufradat}
\kalima{فتى العلا / ابن العلاء}{\AbuAmr.}
\kalima{اليحصبي}{\IbnAmir.}
\kalima{حرمي}{نسبة إلى الحَرَمين: مكة والمدينة.}
\end{mufradat}

\maana[(الأبيات ٥٢–٥٥)] هذه ثماني كلمات، كل كلمة رمز لجماعة:

\begin{qaida}[الرموز الكلمية الثمانية]
\centering
\renewcommand{\arraystretch}{1.5}
\begin{tabular}{c l l}
\toprule
\textbf{الكلمة} & \textbf{القراء} & \textbf{طريقة الحفظ}\\
\midrule
\rzlg{صُحْبَة} & \Hamza، \Kisai، \Shuba & «صحبة» فيها «شعبة»\\
\rzlg{صِحَاب} & \Hamza، \Kisai، \Hafs & «صحاب» مع «حفص»\\
\rzlg{عَمّ} & \Nafi، \IbnAmir & المدني والشامي\\
\rzlg{سَمَا} & \Nafi، \IbnKathir، \AbuAmr & الحجازيان والبصري\\
\rzlg{حَقّ} & \IbnKathir، \AbuAmr & المكي والبصري\\
\rzlg{نَفَر} & \IbnKathir، \AbuAmr، \IbnAmir & «حق» + ابن عامر\\
\rzlg{حِرْمِيّ} & \Nafi، \IbnKathir & أهل الحرمين\\
\rzlg{حِصْن} & \Nafi{} والكوفيون & «ث» + نافع\\
\bottomrule
\end{tabular}
\end{qaida}

\begin{amthila}[أمثلة لاستعمال الرموز من أبيات الشاطبية الآتية]
\begin{itemize}
\item \rzlg{ث}: «وتظاهرون الظاء خفّف \rzg{ثابتا}» \ila{} الكوفيون: \ayah{تَظَٰهَرُونَ} \src{البقرة: ٨٥} بتخفيف الظاء، والباقون \textbf{تَظَّاهرون} بالتشديد.
\item \rzlg{غ}: «عبادُ برفع الدال في عند \rzg{غلغلا}» \ila{} \AbuAmr{} والكوفيون: \ayah{عِبَٰدُ ٱلرَّحۡمَٰنِ} \src{الزخرف: ١٩}، والباقون \textbf{عِنْدَ الرحمن}.
\item \rzlg{صحبة}: «و\rzg{صحبة} يَصرِفْ فتحُ ضمٍّ» \ila{} \Shuba{} و\Hamza{} و\Kisai: \ayah{مَّن يَصۡرِفۡ عَنۡهُ} \src{الأنعام: ١٦} بفتح الياء وكسر الراء، والباقون \textbf{يُصْرَفْ}.
\item \rzlg{صحاب}: «يُضَلُّ بضم الياء مع فتح ضاده \rzg{صحاب}» \ila{} \Hafs{} و\Hamza{} و\Kisai: \ayah{يُضَلُّ بِهِ ٱلَّذِينَ كَفَرُواْ} \src{التوبة: ٣٧}، والباقون \textbf{يَضِلُّ}.
\item \rzlg{حق}: «و\rzg{حقّ} \rz{Asim}{نصيرٍ} كسرُ واوِ مسوِّمين» \ila{} \IbnKathir{} و\AbuAmr{} و\Asim: \ayah{مُسَوِّمِينَ} \src{آل عمران: ١٢٥} بكسر الواو، والباقون بفتحها.
\item \rzlg{حصن}: «وفي المخلَصين الكل \rzg{حصن}» \ila{} \Nafi{} والكوفيون: \ayah{ٱلۡمُخۡلَصِينَ} بفتح اللام حيث وقعت، والباقون بكسرها.
\end{itemize}
\end{amthila}

\bayt{56}{وَمَهْمَا أَتَتْ مِنْ قَبْلُ أَوْ بَعْدُ كِلْمَةٌ}{فَكُنْ عِنْدَ شَرْطِي وَاقْضِ بِالْوَاوِ فَيْصَلَا}
\maana إذا اجتمع رمز كلمي (من الثمانية) مع رمز حرفي، سواء جاءت الكلمة قبل الحرف أو بعده، فالتزم شرطي: كل رمز يدل على ما وُضع له، ولا يتغير شيء، والواو تفصل بين المسائل كما سبق.

\begin{amthila}
\begin{itemize}
\item «\rzg{وحقّ} \rz{Asim}{نصيرٍ}»: \rzg{حق} = \IbnKathir{} و\AbuAmr، و\rz{Asim}{ن} = \Asim{} \ila{} ثلاثتهم.
\item «\rzg{وعمَّ} \rz{Asim}{علا}»: \rzg{عم} = \Nafi{} و\IbnAmir، و\rz{Asim}{ع} = \Hafs.
\end{itemize}
\end{amthila}

\section{القراءة المسكوت عنها: الأضداد}

\bayt{57}{وَمَا كَانَ ذَا ضِدٍّ فَإِنِّي بِضِدِّهِ}{غَنِيٌّ فَزَاحِمْ بِالذَّكَاءِ لِتَفْضُلَا}
\bayt{58}{كَمَدٍّ وَإِثْبَاتٍ وَفَتْحٍ وَمُدْغَمٍ}{وَهَمْزٍ وَنَقْلٍ وَاخْتِلَاسٍ تَحَصَّلَا}
\bayt{59}{وَجَزْمٍ وَتَذْكِيرٍ وَغَيْبٍ وَخِفَّةٍ}{وَجَمْعٍ وَتَنْوِينٍ وَتَحْرِيكٍ اعْمِلَا}
\begin{mufradat}
\kalima{زاحِم بالذكاء}{اجتهد بفطنتك.}
\kalima{لتفضُلا}{لتكون من الفضلاء.}
\end{mufradat}

\maana[(الأبيات ٥٧–٥٩)] هذه قاعدة عظيمة: الناظم يذكر قراءة من رمز لهم فقط، ويسكت عن الباقين. فإذا كانت القراءة المذكورة لها \textbf{ضد}، فقراءة الباقين تكون بذلك الضد.

\begin{amthila}[مثال]
قال في سورة النساء: «وكوفيهم تساءلون \textbf{مخففًا}». فذكر أن الكوفيين يقرؤون بالتخفيف: \ayah{تَسَآءَلُونَ بِهِۦ} \src{النساء: ١}. وسكت عن الباقين، وضد التخفيف التشديد، فالباقون: \textbf{تَسَّاءَلُونَ}.
\end{amthila}

\begin{qaida}[جدول الأضداد]
\centering
\renewcommand{\arraystretch}{1.4}
\begin{tabular}{c c c}
\toprule
\textbf{إذا ذكر} & \textbf{فضدّه للمسكوت عنهم} & \textbf{هل ينعكس؟}\\
\midrule
المد & القصر & نعم\\
الإثبات & الحذف & نعم\\
الفتح & الإمالة & نعم\\
الإدغام & الإظهار & نعم\\
الهمز & تركه & نعم\\
النقل & تركه (إبقاء الحركة) & نعم\\
الاختلاس & إتمام الحركة & نعم\\
الجزم & الرفع & \textbf{لا} (الرفع ضده النصب)\\
التذكير & التأنيث & نعم\\
الغيب & الخطاب & نعم\\
الخِفّة (التخفيف) & الثِّقَل (التشديد) & نعم\\
الجمع & الإفراد & نعم\\
التنوين & تركه & نعم\\
التحريك & الإسكان & نعم\\
\bottomrule
\end{tabular}

\smallskip\raggedleft
«ينعكس» معناه: إذا ذكر القصر فقراءة الآخرين المد، وإذا ذكر المد فقراءتهم القصر. إلا الجزم: ضده الرفع، أما الرفع فضده النصب لا الجزم.
\end{qaida}

\bayt{60}{وَحَيْثُ جَرَى التَّحْرِيكُ غَيْرَ مُقَيَّدٍ}{هُوَ الْفَتْحُ وَالْإِسْكَانُ آخَاهُ مَنْزِلَا}
\begin{mufradat}
\kalima{غير مقيّد}{لم يذكر نوع الحركة.}
\kalima{آخاه منزلًا}{صار ضده المقابل له.}
\end{mufradat}
\maana إذا قال «حرِّك» ولم يذكر الحركة فالمراد \textbf{الفتح}، وضده الإسكان. أما إذا قيّد الحركة (حرِّك بالضم مثلًا) فالمراد ما قيّده، وضده الإسكان أيضًا.
\begin{amthila}
«معًا \textbf{قدر حرِّك} \rz{IbnAmir}{مِن} \rzg{صِحاب}» \ila{} \IbnDhakwan{} و\Hafs{} و\Hamza{} و\Kisai: \ayah{عَلَى ٱلۡمُوسِعِ قَدَرُهُۥ وَعَلَى ٱلۡمُقۡتِرِ قَدَرُهُۥ} \src{البقرة: ٢٣٦} بفتح الدال، والباقون بإسكانها: \textbf{قَدْرُه}.
\end{amthila}

\bayt{61}{وَآخَيْتُ بَيْنَ النُّونِ وَالْيَا وَفَتْحِهِمْ}{وَكَسْرٍ وَبَيْنَ النَّصْبِ وَالْخَفْضِ مُنْزِلَا}
\maana وجعلت هذه الأزواج أضدادًا أيضًا:
\begin{itemize}
\item \textbf{النون والياء}: «وحيث يشاءُ \textbf{نونُ} \rz{IbnKathir}{دارٍ}» \ila{} \IbnKathir: \ayah{يَتَبَوَّأُ مِنۡهَا حَيۡثُ نَشَآءُ} \src{يوسف: ٥٦} بالنون، والباقون \textbf{يشاء} بالياء.
\item \textbf{الفتح والكسر}: «إنّ الدين \textbf{بالفتح} \rz{Kisai}{رُفِّلا}» \ila{} \Kisai: \ayah{أَنَّ ٱلدِّينَ عِندَ ٱللَّهِ ٱلۡإِسۡلَٰمُ} \src{آل عمران: ١٩} بفتح الهمزة، والباقون \textbf{إنّ} بالكسر.
\item \textbf{النصب والخفض}: «\Hamza{} والأرحامِ \textbf{بالخفض} جمّلا» \ila{} \Hamza: \ayah{وَٱلۡأَرۡحَامِ} \src{النساء: ١} بالجر، والباقون \textbf{والأرحامَ} بالنصب.
\end{itemize}
وفرّق الناظم بين ألقاب البناء (الفتح والكسر) وألقاب الإعراب (النصب والخفض) على اصطلاح البصريين.

\bayt{62}{وَحَيْثُ أَقُولُ الضَّمُّ وَالرَّفْعُ سَاكِتًا}{فَغَيْرُهُمْ بِالْفَتْحِ وَالنَّصْبِ أَقْبَلَا}
\maana إذا قلت «الضم» ولم أذكر ضده، فقراءة الباقين \textbf{بالفتح}. وإذا قلت «الرفع» ولم أذكر ضده، فقراءة الباقين \textbf{بالنصب}.
\begin{amthila}
\begin{itemize}
\item «وفي إذ يرون الياءُ \textbf{بالضم} \rz{IbnAmir}{كُلِّلا}» \ila{} \IbnAmir: \ayah{إِذۡ يُرَوۡنَ ٱلۡعَذَابَ} \src{البقرة: ١٦٥} بضم الياء، والباقون \textbf{يَرَوْن} بالفتح.
\item «وحتى يقولَ \textbf{الرفعُ} في اللام \rz{Nafi}{أوّلا}» \ila{} \Nafi: \ayah{حَتَّىٰ يَقُولُ ٱلرَّسُولُ} \src{البقرة: ٢١٤} بالرفع، والباقون \textbf{يقولَ} بالنصب.
\end{itemize}
أما إذا قيّد الضم بأنه ضد الإسكان أو ضد الكسر، فقراءة الباقين بما قيّده.
\end{amthila}

\bayt{63}{وَفِي الرَّفْعِ وَالتَّذْكِيرِ وَالْغَيْبِ جُمْلَةٌ}{عَلَى لَفْظِهَا أَطْلَقْتُ مَنْ قَيَّدَ الْعُلَا}
\begin{mufradat}
\kalima{أطلقتُ}{ذكرتُ الكلمة دون أن أصفها بالرفع أو التذكير أو الغيب.}
\kalima{من قيّد العُلا}{يا من قيّد المعالي (أي حصّلها) — نداء للطالب.}
\end{mufradat}
\maana قد يذكر الناظم الكلمة بلفظها مطلقة دون قيد، فيُفهم أن المراد: \textbf{الرفع} إن كان الخلاف فيها بين الرفع وغيره، و\textbf{التذكير} إن كان بين التذكير والتأنيث، و\textbf{الغيب} إن كان بين الغيب والخطاب.
\begin{amthila}
\begin{itemize}
\item «وأربعُ أولًا \rzg{صحاب}» \ila{} \Hamza{} و\Kisai{} و\Hafs: \ayah{فَشَهَٰدَةُ أَحَدِهِمۡ أَرۡبَعُ شَهَٰدَٰتِۭ} \src{النور: ٦} بالرفع (ولم يقل «بالرفع» بل أطلقها)، والباقون \textbf{أربعَ} بالنصب. وكلمة «أولًا» هنا معناها: الموضع الأول، وليست رمزًا.
\item «بل يؤثرون \rz{AbuAmr}{حُزْ}» \ila{} \AbuAmr: \ayah{بَلۡ يُؤۡثِرُونَ ٱلۡحَيَوٰةَ ٱلدُّنۡيَا} \src{الأعلى: ١٦} بالغيب (الياء)، والباقون \textbf{تؤثرون} بالخطاب.
\end{itemize}
\end{amthila}

\section{موضع الرموز الكلمية، والتصريح بالأسماء}

\bayt{64}{وَقَبْلَ وَبَعْدَ الْحَرْفِ آتِي بِكُلِّ مَا}{رَمَزْتُ بِهِ فِي الْجَمْعِ إِذْ لَيْسَ مُشْكِلَا}
\begin{mufradat}
\kalima{الحرف}{الكلمة القرآنية.}
\kalima{ما رمزت به في الجمع}{الكلمات الثماني (صحبة، صحاب، عم، سما، حق، نفر، حرمي، حصن).}
\end{mufradat}
\maana الرموز الحرفية (أبج...) \textbf{لا تأتي إلا بعد} الكلمة القرآنية. أما الرموز الكلمية الثمانية فقد آتي بها \textbf{قبل} الكلمة القرآنية أو \textbf{بعدها}؛ لأنها لا تُشكل (فهي كلمات معروفة لا تلتبس).
\begin{amthila}
\begin{itemize}
\item بعد الكلمة: «من يرتدِدْ \rzg{عمّ}» \ila{} \Nafi{} و\IbnAmir: \ayah{مَن يَرۡتَدِدۡ مِنكُمۡ} \src{المائدة: ٥٤} بدالين، والباقون \textbf{يرتدَّ}.
\item قبل الكلمة: «و\rzg{صحبة} يَصرِف» كما سبق.
\end{itemize}
وإذا اجتمع رمز حرفي مع رمز كلمي جاز تقديم الحرف تبعًا للكلمة: «و\rzg{حقّ} \rz{Asim}{نصيرٍ} كسرُ واو مسوّمين».
\end{amthila}

\bayt{65}{وَسَوْفَ أُسَمِّي حَيْثُ يَسْمَحُ نَظْمُهُ}{بِهِ مُوضِحًا جِيدًا مُعَمًّا وَمُخْوَلَا}
\begin{mufradat}
\kalima{أسمّي}{أذكر اسم القارئ صريحًا بدل رمزه.}
\kalima{الجِيد}{العنق.}
\kalima{مُعَمًّا ومُخْوَلًا}{كريم الأعمام والأخوال.}
\end{mufradat}
\maana قد أذكر اسم القارئ صريحًا إذا سمح الوزن، فيكون البيان أوضح، كالعنق الجميل الكريم الأعمام والأخوال (أي واضح الحسن). ويأتي الاسم قبل الكلمة القرآنية أو بعدها.
\begin{amthila}
\begin{itemize}
\item قبل الكلمة: «\Hamza{} والأرحامِ بالخفض جمّلا».
\item في أبوابنا: «عليهم إليهم \Hamza{} ولديهمو» (البيت ١١٠)، و«\Qalun{} بتخييره جلا» (البيت ١١١).
\end{itemize}
\end{amthila}

\bayt{66}{وَمَنْ كَانَ ذَا بَابٍ لَهُ فِيهِ مَذْهَبٌ}{فَلَا بُدَّ أَنْ يُسْمَى فَيُدْرَى وَيُعْقَلَا}
\maana إذا انفرد قارئ أو راوٍ بباب كامل لا يشاركه فيه غيره، فلا بد أن أذكره باسمه الصريح لا برمزه.
\begin{amthila}
\begin{itemize}
\item «ودونك الإدغام الكبير وقطبه \AbuAmr{} البصري» (باب الإدغام الكبير).
\item «و\Hamza{} عند الوقف سهّل همزه» (باب وقف حمزة وهشام على الهمز).
\item «وغلّظ \Warsh{} فتح لام» (باب اللامات).
\end{itemize}
\end{amthila}

\begin{khulasa}[خلاصة الاصطلاحات]
\begin{enumerate}
\item الكلمة القرآنية أولًا، ثم رمز الحرف بعدها (من «أبجد»).
\item الواو تفصل بين المسائل.
\item الرموز الكلمية الثمانية تأتي قبل الكلمة أو بعدها.
\item قد يُصرَّح بالاسم، ولا بد منه إذا انفرد القارئ بباب.
\item يُستغنى باللفظ عن القيد إذا كان اللفظ واضحًا.
\item المسكوت عنهم يُقرأ لهم بضد القراءة المذكورة (جدول الأضداد).
\item التحريك المطلق = الفتح. الضم المطلق ضده الفتح، والرفع المطلق ضده النصب.
\item الكلمة المطلقة بلا قيد = الرفع أو التذكير أو الغيب.
\end{enumerate}
\end{khulasa}
      % ٤٤–٦٦  : الرموز والاصطلاحات
% =====================================================================
%  مقدمة الناظم (٤): الأبيات ٦٧–٩٤ — وصف القصيدة وتسميتها، ثم الدعاء والوصايا
% =====================================================================
\chapter{مقدمة الناظم (٤): وصف القصيدة والوصايا}

\section{وصف القصيدة وصلتها بكتاب التيسير}

\bayt{67}{أَهَلَّتْ فَلَبَّتْهَا الْمَعَانِي لُبَابُهَا}{وَصُغْتُ بِهَا مَا سَاغَ عَذْبًا مُسَلْسَلَا}
\begin{mufradat}
\kalima{أهلّت}{رفعت صوتها، أي نادت.}
\kalima{لبّتها}{أجابتها.}
\kalima{لُبابها}{خالصها وأحسنها.}
\kalima{صُغت}{أحكمت النظم (من الصياغة).}
\kalima{ساغ}{سهُل وطاب.}
\kalima{مسلسلًا}{سلِسًا صافيًا.}
\end{mufradat}
\maana نادت القصيدةُ المعانيَ فأجابها خيارها، ونظمتُ فيها لفظًا سهلًا عذبًا يجري على اللسان.

\bayt{68}{وَفِي يُسْرِهَا التَّيْسِيرَ رُمْتُ اخْتِصَارَهُ}{فَأَجْنَتْ بِعَوْنِ اللهِ مِنْهُ مُؤَمَّلَا}
\begin{mufradat}
\kalima{التيسير}{كتاب «التيسير في القراءات السبع» للإمام أبي عمرو الداني (ت ٤٤٤ هـ).}
\kalima{رُمتُ}{قصدتُ.}
\kalima{أجنت}{أثمرت وكثر جناها.}
\kalima{مؤمّلًا}{مرجوًّا.}
\end{mufradat}
\maana قصدت في هذه القصيدة الميسّرة أن أختصر كتاب «التيسير» للداني، فأثمرت بعون الله ما كنت أرجوه منه.

\begin{tanbih}[فائدة]
لهذا يقال: الشاطبية نَظْمُ التيسير. فأصل مادتها من التيسير، مع زيادات قليلة تسمى «زيادات القصيد».
\end{tanbih}

\bayt{69}{وَأَلْفَافُهَا زَادَتْ بِنَشْرِ فَوَائِدٍ}{فَلَفَّتْ حَيَاءً وَجْهَهَا أَنْ تُفَضَّلَا}
\begin{mufradat}
\kalima{ألفافها}{أشجارها الملتفّة، أي أبياتها الكثيرة.}
\kalima{بنشر}{بكثرة.}
\kalima{لفّت وجهها}{سترته.}
\end{mufradat}
\maana زادت القصيدة على التيسير بفوائد كثيرة (كزيادة أحكام، وإشارات لتعليل، ومخارج الحروف)، ومع ذلك غطّت وجهها حياءً أن تُفضَّل على أصلها، أدبًا من الفرع مع الأصل ومن التلميذ مع شيخه.

\bayt{70}{وَسَمَّيْتُهَا حِرْزَ الْأَمَانِي تَيَمُّنًا}{وَوَجْهَ التَّهَانِي فَاهْنِهِ مُتَقَبَّلَا}
\begin{mufradat}
\kalima{الحِرز}{ما يُحفظ فيه الشيء.}
\kalima{الأماني}{جمع أُمنية، وهي ما يُتمنى.}
\kalima{وجه التهاني}{أحسن التهاني (وجه الشيء: أحسنه).}
\kalima{تيمّنًا}{تبرّكًا.}
\kalima{فاهْنِه}{فأعطِه (أصله: فاهنأه).}
\end{mufradat}
\maana سميت القصيدة «\textbf{حِرز الأماني ووجه التهاني}» تبركًا وتفاؤلًا بأن تتحقق فيها أماني طلاب هذا العلم. فيا أيها الطالب أعطها عنايتك متقبلًا لها.

\section{الدعاء}

\bayt{71}{وَنَادَيْتُ اللَّهُمَّ يَا خَيْرَ سَامِعٍ}{أَعِذْنِي مِنَ التَّسْمِيعِ قَوْلًا وَمِفْعَلَا}
\begin{mufradat}
\kalima{أعِذني}{أجِرني واحفظني.}
\kalima{التسميع}{أن يعمل الخير ليسمعه الناس فيمدحوه (الرياء).}
\kalima{مِفْعلًا}{فعلًا.}
\end{mufradat}
\maana يا خير مجيب للدعاء، احفظني من الرياء وطلب السمعة في القول والفعل. وخاف الناظم لمّا أشاد بقصيدته أن يدخله شيء من الرياء فاستعاذ منه.

\bayt{72}{إِلَيْكَ يَدِي مِنْكَ الْأَيَادِي تَمُدُّهَا}{أَجِرْنِي فَلَا أَجْرِي بِجَوْرٍ فَأَخْطَلَا}
\begin{mufradat}
\kalima{الأيادي}{النِّعم.}
\kalima{أجِرني}{احفظني.}
\kalima{الجَور}{الميل عن الحق.}
\kalima{أخطلا}{أقع في الكلام الفاسد.}
\end{mufradat}
\maana يا رب إليك أمدّ يدي، ونِعمُك المتتابعة عليّ هي التي جرّأتني على مدّها. احفظني من الميل عن الحق فأقع في فساد القول.

\bayt{73}{أَمِينَ وَأَمْنًا لِلْأَمِينِ بِسِرِّهَا}{وَإِنْ عَثَرَتْ فَهْوَ الْأَمُونُ تَحَمُّلَا}
\begin{mufradat}
\kalima{أمِين}{لغة في «آمين»: اللهم استجب.}
\kalima{الأمين}{الحافظ الموثوق به.}
\kalima{عثرت}{سقطت وأخطأت.}
\kalima{الأمون}{الناقة القوية التي يؤمَن عِثارها.}
\end{mufradat}
\maana اللهم استجب، وهَبْ أمانًا لمن حفظ أسرار هذه القصيدة ونشرها. وإن وقع فيها خطأ فليتحمّله هذا الأمين ويستره، كما تتحمل الناقة القوية الأثقال.

\bayt{74}{أَقُولُ لِحُرٍّ وَالْمُرُوءَةُ مَرْؤُهَا}{لِإِخْوَتِهِ الْمِرْآةُ ذُو النُّورِ مِكْحَلَا}
\begin{mufradat}
\kalima{الحُرّ}{من لم يستعبده هواه.}
\kalima{المروءة}{كمال الأخلاق.}
\kalima{مرؤها}{صاحبها.}
\kalima{ذو النور}{صاحب الإيمان.}
\kalima{مِكحلًا}{الميل الذي يُكتحل به.}
\end{mufradat}
\maana أقول للكريم: صاحب المروءة لإخوانه كالمرآة يريهم عيوبهم ليصلحوها، وبنور إيمانه يشفي عيوبهم كما يشفي الكحلُ العينَ. وفيه إشارة لحديث: «المؤمن مرآة أخيه» (رواه أبو داود).

\section{وصية القارئ لقصيدته}

\bayt{75}{أَخِي أَيُّهَا الْمُجْتَازُ نَظْمِي بِبَابِهِ}{يُنَادَى عَلَيْهِ كَاسِدَ السُّوقِ أَجْمِلَا}
\bayt{76}{وَظُنَّ بِهِ خَيْرًا وَسَامِحْ نَسِيجَهُ}{بِالْإِغْضَاءِ وَالْحُسْنَى وَإِنْ كَانَ هَلْهَلَا}
\bayt{77}{وَسَلِّمْ لِإِحْدَى الْحُسْنَيَيْنِ إِصَابَةٌ}{وَالُاخْرَى اجْتِهَادٌ رَامَ صَوْبًا فَأَمْحَلَا}
\bayt{78}{وَإِنْ كَانَ خَرْقٌ فَادَّرِكْهُ بِفَضْلَةٍ}{مِنَ الْحِلْمِ وَلْيُصْلِحْهُ مَنْ جَادَ مِقْوَلَا}
\begin{mufradat}
\kalima{المجتاز}{المارّ.}
\kalima{يُنادى عليه}{يُعرض للبيع.}
\kalima{كاسد السوق}{غير مرغوب فيه.}
\kalima{أجمِلا}{قل القول الجميل.}
\kalima{نسيجه}{نظمه (شبّهه بالثوب المنسوج).}
\kalima{الإغضاء}{غضّ الطرف عن العيب.}
\kalima{هلهلًا}{رقيقًا ضعيف النسج.}
\kalima{إصابة}{وصول إلى الصواب.}
\kalima{رام صوبًا فأمحلا}{طلب المطر فأصابه الجدب، أي اجتهد فأخطأ.}
\kalima{خَرْق}{عيب.}
\kalima{ادّرِكه}{تداركه وأصلحه.}
\kalima{من جاد مِقْولًا}{من كان فصيح اللسان متمكنًا من العلم.}
\end{mufradat}
\maana[(الأبيات ٧٥–٧٨)] يا أخي المارّ بقصيدتي وهي تُعرض كأنها بضاعة كاسدة: قل فيها قولًا جميلًا، وأحسن الظن بناظمها، وتجاوز عن عيوبها ولو كان نظمها ضعيفًا كالثوب الرقيق.
وسلّم لي بإحدى الحُسنيين: إما أني أصبت فلي أجران، وإما أني اجتهدت فأخطأت فلي أجر، كما في الحديث: «إذا حكم الحاكم فاجتهد ثم أصاب فله أجران، وإذا حكم فاجتهد ثم أخطأ فله أجر» (متفق عليه).
وإن وجدت عيبًا فتداركه بحلمك، وليصلحه العالِم الفصيح. (وهذا تواضع من الناظم، وإلا فنظمه آية في القوة والإتقان.)

\section{وصايا وآداب}

\bayt{79}{وَقُلْ صَادِقًا لَوْلَا الْوِئَامُ وَرُوحُهُ}{لَطَاحَ الْأَنَامُ الْكُلُّ فِي الْخُلْفِ وَالْقِلَى}
\begin{mufradat}
\kalima{الوئام}{الوفاق والمحبة.}
\kalima{طاح}{هلك.}
\kalima{الأنام}{الناس.}
\kalima{الخُلف}{الاختلاف.}
\kalima{القِلى}{البغض.}
\end{mufradat}
\maana قل صادقًا: لولا الوفاق والمحبة بين الناس لهلكوا جميعًا في الخلاف والبغض. فالوفاق سبب الحياة الطيبة، والخلاف سبب الهلاك.

\bayt{80}{وَعِشْ سَالِمًا صَدْرًا وَعَنْ غِيبَةٍ فَغِبْ}{تُحَضَّرْ حِظَارَ الْقُدْسِ أَنْقَى مُغَسَّلَا}
\begin{mufradat}
\kalima{غِيبة}{ذكر المرء أخاه بما يكره.}
\kalima{فغِب}{ابتعد عن مجالسها.}
\kalima{حِظار القدس}{الجنة (الحِظار: ما يُحاط به، والقدس: الطهر).}
\kalima{أنقى مغسَّلا}{نظيفًا مطهرًا من الذنوب.}
\end{mufradat}
\maana عش سليم القلب من الغش والحقد، وابتعد عن مجالس الغيبة، يُدخلك الله الجنة طاهرًا من الذنوب.

\bayt{81}{وَهَذَا زَمَانُ الصَّبْرِ مَنْ لَكَ بِالَّتِي}{كَقَبْضٍ عَلَى جَمْرٍ فَتَنْجُو مِنَ الْبَلَا}
\maana زماننا زمان الصبر، فمن يأتيك بحالة الصبر التي هي كالقبض على الجمر فتنجو بها من العذاب؟ وفيه إشارة لحديث: «يأتي على الناس زمان الصابر فيهم على دينه كالقابض على الجمر» (رواه الترمذي).

\bayt{82}{وَلَوْ أَنَّ عَيْنًا سَاعَدَتْ لَتَوَكَّفَتْ}{سَحَائِبُهَا بِالدَّمْعِ دِيمًا وَهُطَّلَا}
\bayt{83}{وَلَكِنَّهَا عَنْ قَسْوَةِ الْقَلْبِ قَحْطُهَا}{فَيَا ضَيْعَةَ الْأَعْمَارِ تَمْشِي سَبَهْلَلَا}
\begin{mufradat}
\kalima{ساعدت}{أعانت صاحبها على البكاء.}
\kalima{توكّفت}{قطرت.}
\kalima{سحائبها}{مدامعها (شبّهها بالسحاب).}
\kalima{دِيَمًا}{جمع دِيمة: المطر الدائم.}
\kalima{هُطّلا}{جمع هاطل: المطر المتتابع.}
\kalima{قحطها}{جفافها.}
\kalima{سبهللًا}{فارغة بلا فائدة.}
\end{mufradat}
\maana[(البيتان ٨٢–٨٣)] لو أعانت العينُ صاحبها على البكاء على تقصيره في طاعة الله لسالت دموعها دائمة متتابعة كالمطر، ولكن جفاف العين سببه قسوة القلب من الغفلة. فيا لضياع الأعمار التي تمر فارغة بلا عمل نافع!

\bayt{84}{بِنَفْسِي مَنِ اسْتَهْدَى إِلَى اللهِ وَحْدَهُ}{وَكَانَ لَهُ الْقُرْآنُ شِرْبًا وَمَغْسَلَا}
\bayt{85}{وَطَابَتْ عَلَيْهِ أَرْضُهُ فَتَفَتَّقَتْ}{بِكُلِّ عَبِيرٍ حِينَ أَصْبَحَ مُخْضَلَا}
\begin{mufradat}
\kalima{بنفسي}{أفدي بنفسي.}
\kalima{استهدى}{طلب الهداية.}
\kalima{شِربًا}{نصيبًا من الماء يشربه.}
\kalima{مَغسَلًا}{مكانًا يغتسل فيه، أي يتطهر به من الذنوب.}
\kalima{تفتّقت}{انشقّت.}
\kalima{عَبِير}{أخلاط من الطيب.}
\kalima{مُخضَلًّا}{مبتلًّا.}
\end{mufradat}
\maana[(البيتان ٨٤–٨٥)] أفدي بنفسي من طلب الهداية من الله وحده، وكان القرآن له شرابًا يرويه وماءً يطهره من الذنوب. وطابت له الأرض التي يمشي عليها فانتشر منها الطيب (كناية عن ثناء الناس عليه وكثرة الخير ببركته) حين ابتلّ بنعم الله.

\bayt{86}{فَطُوبَى لَهُ وَالشَّوْقُ يَبْعَثُ هَمَّهُ}{وَزَنْدُ الْأَسَى يَهْتَاجُ فِي الْقَلْبِ مُشْعِلَا}
\bayt{87}{هُوَ الْمُجْتَبَى يَغْدُو عَلَى النَّاسِ كُلِّهِمْ}{قَرِيبًا غَرِيبًا مُسْتَمَالًا مُؤَمَّلَا}
\begin{mufradat}
\kalima{طوبى}{الحالة الطيبة، أو شجرة في الجنة.}
\kalima{همّه}{قصده وإرادته.}
\kalima{زند الأسى}{ما يقدح نار الحزن.}
\kalima{يهتاج}{يثور ويلتهب.}
\kalima{المجتبى}{المختار.}
\kalima{يغدو}{يمرّ ويصبح.}
\kalima{مستمالًا}{يميل الناس إليه.}
\kalima{مؤمَّلًا}{يُرجى عند الشدائد.}
\end{mufradat}
\maana[(البيتان ٨٦–٨٧)] فطوبى لهذا المستهدي، يبعث الشوقُ إرادتَه إلى ما أعدّه الله لأهل طاعته، ويشتعل قلبه حزنًا على ما ضاع من عمره. هو المختار عند الله، يعيش بين الناس قريبًا منهم بتواضعه، غريبًا بندرة حاله، يحبه الناس ويرجون دعاءه.

\bayt{88}{يَعُدُّ جَمِيعَ النَّاسِ مَوْلًى لِأَنَّهُمْ}{عَلَى مَا قَضَاهُ اللهُ يُجْرُونَ أَفْعُلَا}
\bayt{89}{يَرَى نَفْسَهُ بِالذَّمِّ أَوْلَى لِأَنَّهَا}{عَلَى الْمَجْدِ لَمْ تَلْعَقْ مِنَ الصَّبْرِ وَالْأَلَا}
\begin{mufradat}
\kalima{مولى}{سيدًا (على المعنى الأول)، أو عبدًا مملوكًا لله (على المعنى الثاني).}
\kalima{أفعُلا}{أفعالًا.}
\kalima{لم تلعق}{لم تذق.}
\kalima{الصَّبِر}{عصارة شجر مُرّ.}
\kalima{الأَلا}{شجر حسن المنظر مُرّ الطعم.}
\end{mufradat}
\maana[(البيتان ٨٨–٨٩)] هذا المختار يعدّ الناس كلهم سادة تواضعًا فلا يحتقر أحدًا، أو يعدّهم عبيدًا لله لا يملكون له نفعًا ولا ضرًّا فلا يخافهم ولا يتملقهم، لأن أفعالهم تجري بقضاء الله.
ويرى نفسه أولى باللوم من غيره؛ لأنها لم تتحمل مرارة الصبر في طلب المجد عند الله، فيشتغل بعيب نفسه عن عيوب الناس.

\bayt{90}{وَقَدْ قِيلَ كُنْ كَالْكَلْبِ يُقْصِيهِ أَهْلُهُ}{وَمَا يَأْتَلِي فِي نُصْحِهِمْ مُتَبَذِّلَا}
\begin{mufradat}
\kalima{يُقصيه}{يُبعده.}
\kalima{ما يأتلي}{ما يقصّر.}
\kalima{متبذّلًا}{باذلًا كل جهده.}
\end{mufradat}
\maana قيل في المثل: كن كالكلب في الوفاء؛ يطرده أهله ويؤذونه، وهو لا يقصّر في حراستهم ونصحهم. والمقصود: الحث على بذل الجهد في طاعة الله والصبر على البلاء، والإخلاص للناس ولو أساؤوا.

\bayt{91}{لَعَلَّ إِلَهَ الْعَرْشِ يَا إِخْوَتِي يَقِي}{جَمَاعَتَنَا كُلَّ الْمَكَارِهِ هُوَّلَا}
\bayt{92}{وَيَجْعَلُنَا مِمَّنْ يَكُونُ كِتَابُهُ}{شَفِيعًا لَهُمْ إِذْ مَا نَسُوهُ فَيَمْحَلَا}
\begin{mufradat}
\kalima{يقي}{يحفظ.}
\kalima{هُوَّلًا}{جمع هائل: المخيف المفزع.}
\kalima{يَمْحَل}{يشكو ويسعى بصاحبه إلى الله (محَل به: وشى به).}
\end{mufradat}
\maana[(البيتان ٩١–٩٢)] أرجو يا إخوتي أن يحفظنا الله من كل المكاره المفزعة، وأن يجعلنا ممن يكون القرآن شفيعًا لهم يوم القيامة لأنهم لم ينسوه، لا ممن يشكوهم القرآن إلى ربه. وفيه إشارة لحديث: «القرآن شافع مشفَّع، وماحِلٌ مصدَّق».

\bayt{93}{وَبِاللهِ حَوْلِي وَاعْتِصَامِي وَقُوَّتِي}{وَمَا لِيَ إِلَّا سِتْرُهُ مُتَجَلِّلَا}
\bayt{94}{فَيَا رَبِّ أَنْتَ اللهُ حَسْبِي وَعُدَّتِي}{عَلَيْكَ اعْتِمَادِي ضَارِعًا مُتَوَكِّلَا}
\begin{mufradat}
\kalima{حولي}{تحوّلي من حال إلى حال (من المعصية إلى الطاعة).}
\kalima{اعتصامي}{امتناعي من كل ما يَشين.}
\kalima{متجلّلًا}{متغطّيًا.}
\kalima{حسبي}{كافيَّ.}
\kalima{عُدّتي}{ما أُعِدّه لدفع الشدائد.}
\kalima{ضارعًا}{متذللًا.}
\end{mufradat}
\maana[(البيتان ٩٣–٩٤)] تحوّلي من المعصية إلى الطاعة، وامتناعي من الشر، وقوتي على الخير — كل ذلك بالله وحده (وهذا معنى «لا حول ولا قوة إلا بالله»). وليس لي ما أعتمد عليه إلا ستره الذي سترني به في الدنيا، وأرجوه في الآخرة.
يا رب أنت كافيّ في كل أمر، وعُدّتي في كل شدة، وعليك وحدك اعتمادي، متذللًا متوكلًا عليك.

\begin{khulasa}[ختام المقدمة]
بهذا تنتهي مقدمة الشاطبية (٩٤ بيتًا). وأهم ما يجب ضبطه منها قبل الدخول في الأبواب:
\begin{enumerate}
\item أسماء القراء السبعة ورواتهم بالترتيب (الأبيات ٢٥–٤٠).
\item رموز «أبجد» ورموز الجماعة (الأبيات ٤٥–٥٥) — راجع جدول دليل الاستخدام.
\item قواعد الأضداد والإطلاق والتقييد (الأبيات ٥٧–٦٣).
\end{enumerate}
\end{khulasa}
     % ٦٧–٩٤  : تسمية القصيدة والوصايا

% ---------------- أبواب الأصول ----------------
% =====================================================================
%  باب الاستعاذة — الأبيات ٩٥–٩٩
% =====================================================================
\chapter{باب الاستعاذة}

\textbf{الاستعاذة}: طلب العَوذ، أي الالتجاء إلى الله والاحتماء به. والمراد هنا قول القارئ «أعوذ بالله من الشيطان الرجيم» قبل القراءة.
والاستعاذة \textbf{ليست من القرآن} بإجماع العلماء، ولذلك ليس فيها رموز لقراء، بل حكمها عام لجميع القراء. وعدد أبيات الباب خمسة.

\bayt{95}{إِذَا مَا أَرَدْتَ الدَّهْرَ تَقْرَأُ فَاسْتَعِذْ}{جِهَارًا مِنَ الشَّيْطَانِ بِاللهِ مُسْجَلَا}
\bayt{96}{عَلَى مَا أَتَى فِي النَّحْلِ يُسْرًا وَإِنْ تَزِدْ}{لِرَبِّكَ تَنْزِيهًا فَلَسْتَ مُجَهَّلَا}
\begin{mufradat}
\kalima{أردت}{قصدت.}
\kalima{الدهر}{في أي وقت من الزمان.}
\kalima{جِهارًا}{جهرًا، أي بصوت مسموع.}
\kalima{مُسْجَلًا}{مطلقًا: لكل القراء وفي كل القرآن.}
\kalima{على ما أتى في النحل}{بلفظ آية النحل: \ayah{فَإِذَا قَرَأۡتَ ٱلۡقُرۡءَانَ فَٱسۡتَعِذۡ بِٱللَّهِ مِنَ ٱلشَّيۡطَٰنِ ٱلرَّجِيمِ} \src{النحل: ٩٨}.}
\kalima{يُسرًا}{حال كونه سهلًا قليل الكلمات.}
\kalima{تنزيهًا}{تقديسًا وتعظيمًا لله.}
\kalima{مجهَّلًا}{منسوبًا إلى الجهل.}
\end{mufradat}

\maana[(البيتان ٩٥–٩٦)]
إذا أردت أن تقرأ القرآن في أي وقت، ولأي قارئ، ومن أي موضع (أول سورة أو وسطها)، فاستعذ \textbf{جهرًا} قبل القراءة، بلفظ آية النحل:
\begin{center}
\textbf{\large أَعُوذُ بِاللهِ مِنَ الشَّيْطَانِ الرَّجِيمِ}
\end{center}
وهذا هو اللفظ المختار لسهولته وموافقته للقرآن. وإن زدت عليه ألفاظًا فيها تعظيم لله فلا حرج، ولا تُنسب إلى الجهل.

\begin{amthila}[صيغ جائزة فيها زيادة تنزيه]
\begin{itemize}
\item أعوذ بالله \textbf{العظيم} من الشيطان الرجيم.
\item أعوذ بالله \textbf{السميع العليم} من الشيطان الرجيم.
\end{itemize}
\end{amthila}

\begin{qaida}[لماذا قال «إذا ما أردت»؟]
الآية تقول: \ayah{فَإِذَا قَرَأۡتَ ٱلۡقُرۡءَانَ فَٱسۡتَعِذۡ}، وظاهرها أن الاستعاذة بعد القراءة. فبيّن الناظم أن معناها: \textbf{إذا أردتَ} القراءة فاستعذ، أي قبلها. ونظيره: \ayah{إِذَا قُمۡتُمۡ إِلَى ٱلصَّلَوٰةِ فَٱغۡسِلُواْ} \src{المائدة: ٦} أي إذا أردتم القيام إلى الصلاة.
\end{qaida}

\bayt{97}{وَقَدْ ذَكَرُوا لَفْظَ الرَّسُولِ فَلَمْ يَزِدْ}{وَلَوْ صَحَّ هَذَا النَّقْلُ لَمْ يُبْقِ مُجْمَلَا}
\begin{mufradat}
\kalima{ذكروا}{أي بعض العلماء من القراء والمحدثين.}
\kalima{لفظ الرسول}{لفظ استعاذة النبي ﷺ.}
\kalima{مُجمَلًا}{إجمالًا وإطلاقًا في الآية.}
\end{mufradat}
\maana ذكر بعض العلماء أحاديث فيها أن النبي ﷺ لم يزد على لفظ آية النحل في الاستعاذة. ولو صحت هذه الأحاديث لكان اللفظ متعينًا ولم يبق في الآية إجمال، فلا تجوز الزيادة. لكنها أحاديث \textbf{ضعيفة}، وجاءت أحاديث أصح منها فيها زيادة، مثل: «أعوذ بالله السميع العليم من الشيطان الرجيم من همزه ونفخه ونفثه» (رواه أبو داود والترمذي). فبقيت الآية على إطلاقها، فتجوز الزيادة والنقص.

\begin{tanbih}
الأمر بالاستعاذة في الآية \textbf{للندب} (الاستحباب) عند جمهور العلماء من السلف والخلف، لا للوجوب.
\end{tanbih}

\bayt{98}{وَفِيهِ مَقَالٌ فِي الْأُصُولِ فُرُوعُهُ}{فَلَا تَعْدُ مِنْهَا بَاسِقًا وَمُظَلِّلَا}
\begin{mufradat}
\kalima{مقال}{كلام كثير.}
\kalima{الأصول}{كتب أصول الفقه والحديث والقراءات.}
\kalima{فروعه}{تفاصيله (شبّهها بأغصان الشجرة).}
\kalima{لا تَعْدُ}{لا تتجاوز.}
\kalima{باسقًا ومظلِّلًا}{الشجر الطويل الوارف الظل، والمراد: القول القوي المؤيَّد بالدليل.}
\end{mufradat}
\maana في التعوذ كلام طويل في كتب الأصول: هل الأمر للوجوب أو الندب؟ وما درجة الأحاديث فيه؟ وما حكم الجهر والإخفاء والوقف والوصل؟ فارجع إليها ولا تتجاوز القول القوي الذي تؤيده الأدلة.

\bayt{99}{وَإِخْفَاؤُهُ فَصْلٌ أَبَاهُ وُعَاتُنَا}{وَكَمْ مِنْ فَتًى كَالْمَهْدَوِي فِيهِ أَعْمَلَا}
\begin{mufradat}
\kalima{إخفاؤه}{الإسرار بالتعوذ.}
\kalima{فصل}{فرق (بين القرآن وغيره).}
\kalima{أباه}{رفضه.}
\kalima{وُعاتنا}{علماؤنا الحفاظ (جمع واعٍ).}
\kalima{المهدوي}{الإمام أبو العباس أحمد بن عمار المهدوي المقرئ المفسر (ت بعد ٤٣٠ هـ).}
\kalima{أعمل}{أعمل فكره في تصحيحه.}
\end{mufradat}
\maana إخفاء التعوذ (عند من قال به) فيه حكمة: الفرق بين ما هو قرآن وما ليس بقرآن. لكن علماءنا الحفاظ رفضوا الإخفاء وأخذوا بالجهر في جميع القرآن لكل القراء. وكم من عالم كالمهدوي أعمل فكره في تصحيح الإخفاء والإقراء به.

\begin{tanbih}[هل في البيت رمز؟]
حمل بعض الشراح الفاء في «فصل» على أنها رمز \Hamza، والألف في «أباه» على أنها رمز \Nafi، أي أنهما يخفيان التعوذ. والصحيح (كما في الوافي) أنه \textbf{لا رمز في البيت}؛ لأن الباب ليس فيه خلاف بين القراء، والعمل على \textbf{الجهر للجميع}. وقد رُويت عن \Hamza{} و\Nafi{} روايات بالإخفاء من غير طريق الشاطبية.
\end{tanbih}

\begin{qaida}[متى يُستحب إخفاء التعوذ؟ (من الوافي)]
يُستحب الجهر بالتعوذ في عموم الأحوال، ويُستحب الإخفاء في أربعة مواضع:
\begin{enumerate}
\item إذا كان القارئ يقرأ سرًّا، منفردًا أو في مجلس.
\item إذا كان خاليًا (وحده) سواء قرأ سرًّا أو جهرًا.
\item إذا كان في الصلاة (سرية أو جهرية، إمامًا أو مأمومًا أو منفردًا).
\item إذا كان يقرأ وسط جماعة يتدارسون القرآن ولم يكن هو المبتدئ بالقراءة.
\end{enumerate}
وفائدة الجهر: أن يُنصت السامع للقراءة من أولها فلا يفوته منها شيء.
\end{qaida}

\begin{tahrir}
\ithaf{6}{إِذَا مَا أَرَدْتَ الدَّهْرَ تَقْرَأُ فَاسْتَعِذْ}{وَبِالْجَهْرِ عِنْدَ الْكُلِّ فِي الْكُلِّ مُسْجَلَا}
\ithaf{7}{بِشَرْطِ اجْتِمَاعٍ وَابْتِدَاءِ دِرَاسَةٍ}{وَلَا مُخْفِيًا أَوْ فِي الصَّلَاةِ فَفَصِّلَا}
\ithaf{8}{وَوَقْفٌ عَلَيْهِ ثُمَّ وَصْلٌ بِأَرْبَعٍ}{لَهُمْ وَاسْتَعِذْ نَدْبًا أَوَ اوْجِبْ وَوَهِّلَا}
\medskip
\textbf{الشرح:}
\begin{itemize}
\item الجهر بالتعوذ لكل القراء في كل القرآن، \textbf{بشرط}: أن يكون معه من يسمعه، وأن يكون هو المبتدئ بالقراءة في مجلس الدرس، وألا يكون يقرأ سرًّا، وألا يكون في الصلاة. فإن اختلّ شرط أخفاه. (وهذا موافق للمواضع الأربعة المذكورة في الوافي.)
\item يجوز \textbf{الوقف} على التعوذ، ويجوز \textbf{وصله} بما بعده، لكل القراء.
\item \textbf{بأربع}: إذا اجتمعت الاستعاذة والبسملة وأول السورة ففيها \textbf{أربعة أوجه} لكل القراء (انظرها في الجدول أدناه).
\item الاستعاذة \textbf{مندوبة}، وقيل واجبة، وهذا القول الثاني ضعيف (معنى «وهِّلا»: اعتبره وهمًا).
\end{itemize}
\end{tahrir}

\begin{khulasa}[الأوجه الأربعة عند اجتماع الاستعاذة والبسملة وأول السورة]
\centering
\renewcommand{\arraystretch}{1.5}
\begin{tabular}{c l l}
\toprule
\textbf{الوجه} & \textbf{الاسم} & \textbf{الصورة}\\
\midrule
١ & قطع الجميع & أعوذ بالله... \textbf{(وقف)} بسم الله... \textbf{(وقف)} الحمد لله...\\
٢ & قطع الأول ووصل الثاني بالثالث & أعوذ بالله... \textbf{(وقف)} بسم الله... الحمد لله...\\
٣ & وصل الأول بالثاني وقطع الثالث & أعوذ بالله... بسم الله... \textbf{(وقف)} الحمد لله...\\
٤ & وصل الجميع & أعوذ بالله... بسم الله... الحمد لله...\\
\bottomrule
\end{tabular}
\end{khulasa}

\begin{tanbih}[تتمة: إذا قطع القارئ قراءته]
لو قطع القارئ قراءته لأمر طارئ قهري (كعطاس أو تنحنح) أو لكلام يتعلق بالقراءة (كسؤال عن شيء فيها)، فإنه \textbf{لا يعيد} التعوذ.
أما لو قطعها إعراضًا عنها، أو لكلام لا يتعلق بها ولو ردًّا لسلام، فإنه \textbf{يستأنف} التعوذ.
\end{tanbih}
             % ٩٥–٩٩
% =====================================================================
%  باب البسملة — الأبيات ١٠٠–١٠٧
% =====================================================================
\chapter{باب البسملة}

\textbf{البسملة}: قول «بسم الله الرحمن الرحيم»، مصدر «بَسْمَلَ» مثل: هَلَّلَ (قال لا إله إلا الله)، وحَمْدَلَ (قال الحمد لله)، وحَوْقَلَ (قال لا حول ولا قوة إلا بالله).

يتكلم الباب عن ثلاثة مواضع:
\begin{enumerate}
\item \textbf{بين السورتين} (آخر سورة مع أول التي بعدها) — وفيه خلاف القراء. \hfill (الأبيات ١٠٠–١٠٤)
\item \textbf{أول سورة براءة}. \hfill (البيت ١٠٥)
\item \textbf{الابتداء بأول السورة أو بوسطها}، وحكم الوقف على البسملة. \hfill (البيتان ١٠٦–١٠٧)
\end{enumerate}

\section{البسملة بين السورتين}

\bayt{100}{وَبَسْمَلَ بَيْنَ السُّورَتَيْنِ \rz{Nafi}{بِسُنَّةٍ}}{\rz{Kisai}{رِجَالٌ} \rz{Asim}{نَمَوْهَا} \rz{IbnKathir}{دِرْيَةً} وَتَحَمُّلَا}
\begin{mufradat}
\kalima{بسنّة}{بطريقة النبي ﷺ. والباء رمز \Qalun.}
\kalima{رجال}{الراء رمز \Kisai.}
\kalima{نَمَوها}{رفعوها ونقلوها (إلى النبي ﷺ). والنون رمز \Asim.}
\kalima{دِرية}{علمًا ومعرفة. والدال رمز \IbnKathir.}
\kalima{تحمّلا}{نقلًا عن الشيوخ (رواية).}
\end{mufradat}
\maana قرأ \Qalun{} و\Kisai{} و\Asim{} و\IbnKathir{} بإثبات البسملة بين كل سورتين، متمسكين في ذلك بسنة نقلوها عن النبي ﷺ، جامعين بين الدراية والرواية.
ومن هذه السنة: أن النبي ﷺ كان لا يعرف انقضاء السورة حتى ينزل عليه «بسم الله الرحمن الرحيم»، وكتابة الصحابة لها في المصاحف العثمانية.

\bayt{101}{وَوَصْلُكَ بَيْنَ السُّورَتَيْنِ \rz{Hamza}{فَصَاحَةٌ}}{وَصِلْ وَاسْكُتَنْ \rz{IbnAmir}{كُلٌّ} \rz{Nafi}{جَلَايَاهُ} \rz{AbuAmr}{حَصَّلَا}}
\begin{mufradat}
\kalima{فصاحة}{بيان. والفاء رمز \Hamza.}
\kalima{كلٌّ}{الكاف رمز \IbnAmir.}
\kalima{جلاياه}{الأمور الواضحة (جمع جَلِيّة). والجيم رمز \Warsh.}
\kalima{حصّلا}{الحاء رمز \AbuAmr.}
\end{mufradat}
\maana
\textbf{أولًا}: \Hamza{} يصل آخر السورة بأول التي بعدها \textbf{بدون بسملة}. وفي قوله «فصاحة» إشارة إلى حكمة الوصل: أنه يُظهر إعراب آخر السورة، كما في آخر التوبة مع أول يونس \ayah{وَهُوَ رَبُّ ٱلۡعَرۡشِ ٱلۡعَظِيمِ ۝ الٓرۚ}، ويُظهر حال همزة الوصل كآخر العاديات مع أول القارعة.

\textbf{ثانيًا}: \IbnAmir{} و\Warsh{} و\AbuAmr{} لكل منهم وجهان بين السورتين (من غير بسملة): \textbf{الوصل} و\textbf{السكت}.

\begin{qaida}[ما هو السكت بين السورتين؟]
\textbf{السكت}: قطع الصوت على آخر السورة زمنًا يسيرًا \textbf{دون تنفس}، بنية مواصلة القراءة، ثم تبدأ السورة التالية بلا بسملة. (وهو أقصر من الوقف الذي يكون معه تنفس.)
\end{qaida}

\begin{tanbih}[وصل حمزة مع سكت خلف]
\Khalaf{} عن \Hamza{} يسكت على الساكن قبل الهمز (كما سيأتي في بابه)، فإذا وصل آخر الضحى \ayah{فَحَدِّثۡ} بأول الشرح \ayah{أَلَمۡ نَشۡرَحۡ} سكت على الثاء. وهذا لا يُخرجه عن كونه \textbf{وصلًا}؛ لأن سكت خلف هذا لا يكون إلا في حال الوصل.
\end{tanbih}

\bayt{102}{وَلَا نَصَّ كَلَّا حُبَّ وَجْهٌ ذَكَرْتُهُ}{وَفِيهَا خِلَافٌ جِيدُهُ وَاضِحُ الطُّلَى}
\begin{mufradat}
\kalima{كلّا}{حرف ردع وزجر.}
\kalima{حُبّ وجهٌ}{ما أحبّ هذا الوجه! (وجهٌ: يعني الوصل والسكت).}
\kalima{جيده}{عنقه.}
\kalima{الطُّلى}{صفحات العنق، و«واضح الطلى» كناية عن الشهرة والوضوح.}
\end{mufradat}
\maana لم يرد نصٌّ عن \IbnAmir{} و\Warsh{} و\AbuAmr{} بالوصل ولا بالسكت، وإنما التخيير بين الوجهين اختيار من أهل الأداء واستحباب من شيوخ الإقراء. وكأن الناظم يزجر («كلّا») من يظن أن هناك نصًّا عنهم، ثم يمدح هذا الوجه («حُبّ وجهٌ»).
ثم قال: \textbf{وفيها} (أي في البسملة) \textbf{خلافٌ} مشهور واضح عن هؤلاء الثلاثة: فيُروى لهم أيضًا البسملة بين السورتين.

\begin{tanbih}[ليس في البيت رمز]
الكاف في «كلّا» والحاء في «حبّ» والجيم في «جيده» \textbf{ليست رموزًا}؛ لأن البيت كله في حكم الثلاثة المذكورين قبله. وهذا ما عليه المحققون.
\end{tanbih}

\begin{khulasa}[حاصل الأبيات ١٠٠–١٠٢: الأوجه بين السورتين]
\centering
\renewcommand{\arraystretch}{1.5}
\begin{tabular}{l l}
\toprule
\textbf{القراء} & \textbf{الأوجه بين السورتين}\\
\midrule
\Qalun، \IbnKathir، \Asim، \Kisai & البسملة فقط\\
\Warsh، \AbuAmr، \IbnAmir & البسملة، والوصل، والسكت\\
\Hamza & الوصل فقط (بلا بسملة)\\
\bottomrule
\end{tabular}
\end{khulasa}

\begin{tahrir}
\ithaf{9}{وَفِيهَا خِلَافٌ جِيدُهُ وَاضِحُ الطُّلَا}{وَذُو الْخُلْفِ لِلْبَصْرِيِّ وَشَامٍ تَنَقَّلَا}
\medskip
قيّد صاحب الإتحاف قول الشاطبي «وفيها خلاف» بأن وجه البسملة إنما نُقل لـ\AbiAmr{} البصري و\IbnAmir{} الشامي. ومقتضى ذلك أن \Warsh{} من طريق الشاطبية له بين السورتين الوصل والسكت فقط.
وهذه من المسائل التي يختلف فيها أهل التحرير، فخذ فيها بما يقرئك به شيخك.
\end{tahrir}

\section{البسملة بين السور الأربع الزُّهر}

\bayt{103}{وَسَكْتُهُمُ الْمُخْتَارُ دُونَ تَنَفُّسٍ}{وَبَعْضُهُمُ فِي الْأَرْبَعِ الزُّهْرِ بَسْمَلَا}
\bayt{104}{لَهُمْ دُونَ نَصٍّ وَهْوَ فِيهِنَّ سَاكِتٌ}{\qname{Hamza}{لِحَمْزَةَ} فَافْهَمْهُ وَلَيْسَ مُخَذَّلَا}
\begin{mufradat}
\kalima{سكتهم}{أي سكت الثلاثة: \Warsh{} و\AbuAmr{} و\IbnAmir.}
\kalima{المختار}{المقدَّم على الوصل.}
\kalima{الأربع الزُّهر}{أربع سور مشهورة: القيامة، والمطففين، والبلد، والهُمَزة. (الزُّهر: المشرقة، كناية عن شهرتها.)}
\kalima{مخذَّلًا}{متروك النصرة، أي ضعيفًا.}
\end{mufradat}
\maana[(البيتان ١٠٣–١٠٤)]
\textbf{أولًا}: السكت للثلاثة (\Warsh{} و\AbuAmr{} و\IbnAmir) هو الوجه المختار المقدَّم على الوصل، لأن فيه تنبيهًا على انتهاء السورة. ويكون بلا تنفس.

\textbf{ثانيًا}: بعض أهل الأداء استحب للقراء مذهبًا خاصًا في أربع سور هي: \textbf{القيامة، المطففين، البلد، الهمزة}، وذلك لبشاعة اللفظ عند وصلها بما قبلها (مثل وصل \ayah{وَأَهۡلُ ٱلۡمَغۡفِرَةِ} بـ\ayah{لَآ أُقۡسِمُ}، أو \ayah{وَٱلۡأَمۡرُ يَوۡمَئِذࣲ لِّلَّهِ} بـ\ayah{وَيۡلࣱ لِّلۡمُطَفِّفِينَ}):
\begin{itemize}
\item من كان يسكت في غيرها (\Warsh{} و\AbuAmr{} و\IbnAmir) \ila{} \textbf{يبسمل} بين هذه الأربع وما قبلها.
\item من كان يصل في غيرها (هؤلاء الثلاثة في وجه الوصل، و\Hamza) \ila{} \textbf{يسكت} بين هذه الأربع وما قبلها.
\end{itemize}
وهذا كله بلا نص عن القراء، بل هو اختيار. والناظم يقول: «فافهمه وليس مخذَّلا» أي هذا المذهب ليس ضعيفًا. لكن \textbf{المحققون على عدم التفريق} بين هذه السور وغيرها، وهو المذهب المختار الذي عليه العمل.

\begin{amthila}[المواضع الأربعة]
\begin{itemize}
\item بين المدثر و\textbf{القيامة}: \ayah{هُوَ أَهۡلُ ٱلتَّقۡوَىٰ وَأَهۡلُ ٱلۡمَغۡفِرَةِ ۝ لَآ أُقۡسِمُ بِيَوۡمِ ٱلۡقِيَٰمَةِ}.
\item بين الانفطار و\textbf{المطففين}: \ayah{وَٱلۡأَمۡرُ يَوۡمَئِذࣲ لِّلَّهِ ۝ وَيۡلࣱ لِّلۡمُطَفِّفِينَ}.
\item بين الفجر و\textbf{البلد}: \ayah{وَٱدۡخُلِي جَنَّتِي ۝ لَآ أُقۡسِمُ بِهَٰذَا ٱلۡبَلَدِ}.
\item بين العصر و\textbf{الهمزة}: \ayah{وَتَوَاصَوۡاْ بِٱلصَّبۡرِ ۝ وَيۡلࣱ لِّكُلِّ هُمَزَةࣲ}.
\end{itemize}
\end{amthila}

\begin{tahrir}
\ithaf{10}{وَبَسْمِلْ بِزُهْرٍ إِنْ تُبَسْمِلْ بِغَيْرِهَا}{وَإِنْ تَسْكُتُ اسْكُتْ بَعْدَ مَا أَنْ تُبَسْمِلَا}
\ithaf{11}{وَإِنْ تَصِلَنْ فَاسْكُتْ بِهَا ثُمَّ صِلْ وَإِنْ}{بَدَأْتَ بِهَا بَسْمِلْ بِهَا وَبِمَا تَلَا}
\ithaf{12}{فَبَسْمِلْ كَذَا اسْكُتْ ثُمَّ إِنْ تَسْكُتَا بِهَا}{فَفِي غَيْرِهَا اسْكُتْ صِلْ وَإِنْ تَصِلًا صِلَا}
\medskip
\textbf{خلاصة ما يتضح من هذه الأبيات} (لمن أخذ بمذهب التفريق في الزُّهر):
\begin{itemize}
\item إن كنت تبسمل في غير الزهر \ila{} فبسمل في الزهر أيضًا.
\item إن كنت تسكت في غيرها \ila{} ففي الزهر البسملة، ويجوز السكت أيضًا.
\item إن كنت تصل في غيرها \ila{} ففي الزهر السكت، ويجوز الوصل أيضًا.
\end{itemize}
وفي بقية البيتين ١١–١٢ تفصيلات دقيقة في ترتيب الأوجه عند الجمع، تُتلقّى من الشيخ عند الأداء.
\end{tahrir}

\section{أول سورة براءة}

\bayt{105}{وَمَهْمَا تَصِلْهَا أَوْ بَدَأْتَ بَرَاءَةً}{لِتَنْزِيلِهَا بِالسَّيْفِ لَسْتَ مُبَسْمِلَا}
\begin{mufradat}
\kalima{تصلها}{تصل سورة براءة (التوبة) بالأنفال قبلها.}
\kalima{بالسيف}{بالقتال والتهديد ونبذ العهود.}
\end{mufradat}
\maana سورة براءة (التوبة) \textbf{لا بسملة في أولها لأحد من القراء}، سواء وصلتها بالأنفال قبلها أو ابتدأت القراءة بها. والعلة: أنها نزلت بالسيف، أي بالأمر بالقتال ونبذ العهود والوعيد. وقد نُقل عن علي رضي الله عنه: «بسم الله أمان، وبراءة ليس فيها أمان؛ نزلت بالسيف».

\begin{tahrir}
\ithaf{13}{وَلِلْكُلِّ قِفْ صِلْ فِي عَلِيمٍ بَرَاءَةٌ}{أَوِ اسْكُتْ وَبَيْنَ النَّاسِ وَالْحَمْدِ بَسْمِلَا}
\medskip
\begin{itemize}
\item بين آخر الأنفال \ayah{إِنَّ ٱللَّهَ بِكُلِّ شَيۡءٍ عَلِيمُۢ} وأول براءة: لكل القراء ثلاثة أوجه: \textbf{الوقف، والسكت، والوصل} (كلها بلا بسملة).
\item بين آخر سورة الناس وأول الفاتحة: \textbf{البسملة لكل القراء} (حتى من مذهبه الوصل أو السكت)؛ لأن الفاتحة مبتدأ بها حكمًا، إذ ليس قبلها شيء في الحقيقة.
\end{itemize}
\end{tahrir}

\section{الابتداء بأول السورة وبوسطها}

\bayt{106}{وَلَا بُدَّ مِنْهَا فِي ابْتِدَائِكَ سُورَةً}{سِوَاهَا وَفِي الْأَجْزَاءِ خُيِّرَ مَنْ تَلَا}
\begin{mufradat}
\kalima{منها}{من البسملة.}
\kalima{سواها}{غير براءة.}
\kalima{الأجزاء}{ما بعد أول السورة ولو بآية أو كلمة (أواسط السور).}
\kalima{خُيِّر من تلا}{خُيِّر القارئ.}
\end{mufradat}
\maana
\begin{itemize}
\item إذا \textbf{ابتدأت} القراءة بأول أي سورة غير براءة \ila{} \textbf{لا بد} من البسملة لجميع القراء، حتى \Hamza{} ومن له السكت أو الوصل؛ لأن مذاهبهم تلك إنما هي \textbf{بين} السورتين.
\item إذا ابتدأت من \textbf{وسط السورة} (ولو من الآية الثانية) \ila{} القارئ مخيّر بين البسملة وتركها، لكل القراء.
\end{itemize}

\begin{tanbih}[الابتداء بوسط براءة]
لا فرق في هذا التخيير بين أجزاء براءة وأجزاء غيرها على الصحيح، وإن منع بعضهم البسملة في أجزاء براءة إلحاقًا لها بأولها.
\end{tanbih}

\bayt{107}{وَمَهْمَا تَصِلْهَا مَعْ أَوَاخِرِ سُورَةٍ}{فَلَا تَقِفَنَّ الدَّهْرَ فِيهَا فَتَثْقُلَا}
\begin{mufradat}
\kalima{تصلها}{تصل البسملة.}
\kalima{فيها}{عليها.}
\kalima{تثقلا}{تصير مستثقَلًا عند أئمة القراءة.}
\end{mufradat}
\maana إذا وصلت البسملة بآخر السورة فلا تقف عليها أبدًا، بل يجب وصلها بأول السورة التالية؛ لأن البسملة شُرعت لأول السورة لا لآخرها، فإذا وقفت عليها بعد وصلها بآخر السورة أوهمت أنها من آخر السورة السابقة.

\begin{khulasa}[الأوجه بين السورتين لمن يبسمل]
\centering
\renewcommand{\arraystretch}{1.5}
\begin{tabular}{c l c}
\toprule
\textbf{الوجه} & \textbf{الصورة} & \textbf{الحكم}\\
\midrule
١ & قطع الجميع: آخر السورة \textbf{(وقف)} البسملة \textbf{(وقف)} أول السورة & جائز\\
٢ & قطع الأول ووصل الثاني بالثالث & جائز\\
٣ & وصل الجميع & جائز\\
٤ & وصل آخر السورة بالبسملة ثم الوقف عليها & \textbf{ممنوع}\\
\bottomrule
\end{tabular}
\end{khulasa}

\section{خلاصة الباب}
\begin{khulasa}[ما لكل قارئ بين كل سورتين (ما عدا الأنفال وبراءة)]
\centering
\renewcommand{\arraystretch}{1.5}
\begin{tabular}{l c l}
\toprule
\textbf{القراء} & \textbf{عدد الأوجه} & \textbf{تفصيلها}\\
\midrule
\Qalun، \IbnKathir، \Asim، \Kisai & ٣ & أوجه البسملة الثلاثة الجائزة\\
\Warsh، \AbuAmr، \IbnAmir & ٥ & أوجه البسملة الثلاثة + الوصل + السكت\\
\Hamza & ١ & الوصل بلا بسملة\\
\midrule
الجميع: بين الأنفال وبراءة & ٣ & الوقف، والسكت، والوصل (بلا بسملة)\\
الجميع: بين الناس والفاتحة & — & البسملة\\
الجميع: الابتداء بأول سورة & — & البسملة (إلا براءة)\\
الجميع: الابتداء بوسط سورة & — & التخيير\\
\bottomrule
\end{tabular}
\end{khulasa}
              % ١٠٠–١٠٧
% =====================================================================
%  سورة أم القرآن — الأبيات ١٠٨–١١٥
% =====================================================================
\chapter{سورة أم القرآن}

يذكر الناظم في هذا الباب خلاف القراء في سورة الفاتحة (وتسمى أم القرآن)، ثم يجعل بعض كلماتها أصولًا مطّردة في القرآن كله: كلمة \textbf{الصراط}، وهاء \textbf{عليهم وإليهم ولديهم}، و\textbf{ميم الجمع}.
وعدد أبياته ثمانية.

\begin{tanbih}[قبل أن تبدأ: كيف تقرأ رموز هذا الباب]
الرمز يأتي \textbf{بعد} الكلمة القرآنية. فانظر إلى أول حرف من الكلمة التي تأتي بعدها: إن كان من حروف «أبجد هوز...» فهو رمز قارئ (راجع جدول الرموز في دليل الاستخدام).
وفي الأبيات ظلّلنا الكلمة القرآنية المختلف فيها بالأصفر \qw{هكذا}، ولوّنّا كلمة الرمز بلون القارئ.
\end{tanbih}

% ---------------------------------------------------------------------
\section{﴿مالك يوم الدين﴾ و﴿الصراط﴾}
% ---------------------------------------------------------------------
\bayt{108}{\qw{وَمَالِكِ} يَوْمِ الدِّينِ \rz{Kisai}{رَاوِيهِ} \rz{Asim}{نَاصِرٌ}}{وَعِنْدَ \qw{سِرَاطٍ} \qw{وَالسِّرَاطَ} \qname{IbnKathir}{لِ قُنْبُلَا}}

\begin{mufradat}
\kalima{رَاوِيهِ}{الراء رمز \Kisai. ومعنى الكلمة: ناقله.}
\kalima{نَاصِرٌ}{النون رمز \Asim. ومعناها: مؤيِّده.}
\kalima{لِ قُنْبُلَا}{«لِ» فعل أمر من «وَلِيَ يَلِي» أي: اتْبَعْ قنبلًا. وليست اللام هنا رمز \Hisham{}؛ لأنه صرّح بعدها باسم \Qunbul.}
\end{mufradat}

\maana
قرأ \Kisai{} و\Asim{} \ayah{مَالِكِ يَوْمِ الدِّينِ} بإثبات الألف بعد الميم، ولم يحتج الناظم إلى أن يقول «بالمد» لأن لفظه بالكلمة «مالك» يدل على القراءة (وهذا معنى قوله في المقدمة: «وباللفظ أستغني عن القيد إن جلا»).
فتكون قراءة الباقين بحذف الألف: \ayah{مَلِكِ يَوْمِ الدِّينِ}.

ثم قال: اتبع \Qunbul{} في قراءة \textbf{سراط} و\textbf{السراط} بالسين حيث وقع في القرآن.

\begin{qiraat}
\qira{\Asim{} و\Kisai}{\ayah{مَالِكِ} بالألف}
\qira{\Baqun{} (\Nafi، \IbnKathir، \AbuAmr، \IbnAmir، \Hamza)}{\ayah{مَلِكِ} بلا ألف}
\qira{\Qunbul}{\ayah{السِّرَاطَ} و\ayah{سِرَاطَ} بالسين في كل القرآن}
\end{qiraat}

\begin{qaida}[قاعدة كلمة «صراط» عند قنبل]
تقرأ لقنبل بالسين في كل موضع، سواء كانت الكلمة:
\begin{itemize}
\item \textbf{معرّفة بأل}: \ayah{ٱهۡدِنَا ٱلسِّرَٰطَ ٱلۡمُسۡتَقِيمَ} \src{الفاتحة: ٦}
\item \textbf{نكرة}: \ayah{وَإِنَّكَ لَتَهۡدِيٓ إِلَىٰ سِرَٰطࣲ مُّسۡتَقِيمࣲ} \src{الشورى: ٥٢}
\item \textbf{مضافة}: \ayah{سِرَٰطَ ٱلَّذِينَ أَنۡعَمۡتَ عَلَيۡهِمۡ} \src{الفاتحة: ٧}، \ayah{وَأَنَّ هَٰذَا سِرَٰطِي مُسۡتَقِيمࣰا} \src{الأنعام: ١٥٣}
\end{itemize}
وكلمة «حيث أتى» في أول البيت التالي هي التي أفادت هذا التعميم.
\end{qaida}

% ---------------------------------------------------------------------
\bayt{109}{بِحَيْثُ أَتَى وَالصَّادَ زَايًا أَشِمَّهَا}{لَدَى \qname{Hamza}{خَلَفٍ} وَاشْمِمْ \qname{Hamza}{لِخَلَّادٍ} الْاَوَّلَا}

\begin{mufradat}
\kalima{بحيث أتى}{أي في كل موضع جاءت فيه الكلمة في القرآن.}
\kalima{أَشِمَّها زايًا}{الإشمام هنا: أن تخلط صوت الصاد بصوت الزاي.}
\kalima{الأوّلا}{الموضع الأول فقط، وهو \ayah{ٱهۡدِنَا ٱلصِّرَٰطَ ٱلۡمُسۡتَقِيمَ} في الفاتحة.}
\end{mufradat}

\maana
قرأ \Khalaf{} بإشمام الصاد صوت الزاي في «صراط» و«الصراط» حيث وقع، وقرأ \Khallad{} بالإشمام في الموضع الأول فقط وهو \ayah{ٱهۡدِنَا ٱلصِّرَٰطَ ٱلۡمُسۡتَقِيمَ}، وما سواه بالصاد الخالصة.
وقرأ الباقون (ومعهم \Khallad{} في غير الموضع الأول) بالصاد الخالصة.

\begin{qiraat}
\qira{\Qunbul}{بالسين: \ayah{ٱلسِّرَٰطَ}}
\qira{\Khalaf}{بالإشمام (صاد مخلوطة بالزاي) في كل القرآن}
\qira{\Khallad}{بالإشمام في \ayah{ٱهۡدِنَا ٱلصِّرَٰطَ} فقط، والباقي بالصاد}
\qira{\Baqun}{بالصاد الخالصة}
\end{qiraat}

\begin{qaida}[كيف يُنطق الإشمام هنا؟]
تخلط لفظ الصاد بلفظ الزاي فيتولد حرف ليس صادًا خالصًا ولا زايًا خالصًا، لكن صوت الصاد هو الغالب.
وتقريبه: أن تنطق الصاد قريبًا من نطق العامّة للظاء (كما قال صاحب الوافي).
\end{qaida}

\begin{tanbih}[تنبيه من الوافي]
ذكر الشيخ القاضي في حاشية الوافي أن \Khallad{} في الموضع الأول (\ayah{ٱهۡدِنَا ٱلصِّرَٰطَ}) له وجهان: الإشمام والصاد الخالصة، وأن اقتصار الناظم على الإشمام فيه قصور. فاعرض ذلك على شيخك عند الأداء.
\end{tanbih}

% ---------------------------------------------------------------------
\section{هاء ﴿عليهم﴾ و﴿إليهم﴾ و﴿لديهم﴾}
% ---------------------------------------------------------------------
\bayt{110}{\qw{عَلَيْهِمْ} \qw{إِلَيْهِمْ} \qname{Hamza}{حَمْزَةٌ} \qw{وَلَدَيْهِمُو}}{جَمِيعًا بِضَمِّ الْهَاءِ وَقْفًا وَمَوْصِلَا}

\begin{mufradat}
\kalima{وقفًا وموصلًا}{أي في حال الوقف وفي حال الوصل.}
\end{mufradat}

\maana
قرأ \Hamza{} هذه الكلمات الثلاث فقط: \textbf{عليهم، إليهم، لديهم} بضم الهاء في كل القرآن، وقفًا ووصلًا، سواء جاء بعدها متحرك أو ساكن.
وقرأ الباقون بكسر الهاء (وأُخذ الكسر من لفظ الناظم بالكلمات مكسورة).

\begin{qiraat}
\qira{\Hamza}{\ayah{عَلَيۡهُمۡ}، \ayah{إِلَيۡهُمۡ}، \ayah{لَدَيۡهُمۡ} بضم الهاء}
\qira{\Baqun}{\ayah{عَلَيۡهِمۡ}، \ayah{إِلَيۡهِمۡ}، \ayah{لَدَيۡهِمۡ} بكسر الهاء}
\end{qiraat}

\begin{amthila}
\begin{itemize}
\item بعدها متحرك: \ayah{غَيۡرِ ٱلۡمَغۡضُوبِ عَلَيۡهِمۡ وَلَا ٱلضَّآلِّينَ} \src{الفاتحة: ٧} \ila{} \Hamza: \textbf{عليهُمْ}.
\item \ayah{يُحِبُّونَ مَنۡ هَاجَرَ إِلَيۡهِمۡ} \src{الحشر: ٩}، \ayah{وَمَا كُنتَ لَدَيۡهِمۡ إِذۡ يُلۡقُونَ} \src{آل عمران: ٤٤}.
\item بعدها ساكن: \ayah{حَقَّ عَلَيۡهِمُ ٱلۡقَوۡلُ} \src{القصص: ٦٣}، \ayah{إِذۡ أَرۡسَلۡنَآ إِلَيۡهِمُ ٱثۡنَيۡنِ} \src{يس: ١٤} \ila{} \Hamza{} يضم الهاء هنا أيضًا.
\end{itemize}
\end{amthila}

\begin{tanbih}
الحكم خاص بهذه الكلمات الثلاث بلفظها. أما نحو \textbf{فيهم، بهم، يزكيهم} فلا يضم \Hamza{} هاءها، بل يكسرها كالباقين (وسيأتي حكمها عند ساكن في البيت ١١٤).
\end{tanbih}

% ---------------------------------------------------------------------
\section{ميم الجمع}
% ---------------------------------------------------------------------
\begin{qaida}[ما هي ميم الجمع؟ وما معنى صلتها؟]
\textbf{ميم الجمع}: الميم الزائدة الدالة على جمع الذكور، مثل: عليه\textbf{م}، أنت\textbf{م}، ك\textbf{م}.\\
\textbf{الصلة}: أن تضم الميم وتمدها بواو لفظية مقدار حركتين وصلًا: \textbf{عليهمُو}، \textbf{أنتمُو}.\\
\textbf{الإسكان}: أن تنطقها ساكنة: \textbf{عليهمْ}.\\
والخلاف كله \textbf{في الوصل فقط}. أما عند الوقف عليها فالجميع يسكّنها.
\end{qaida}

\bayt{111}{وَصِلْ ضَمَّ مِيمِ الْجَمْعِ قَبْلَ مُحَرَّكٍ}{\rz{IbnKathir}{دِرَاكًا} \qname{Nafi}{وَقَالُونٌ} بِتَخْيِيرِهِ جَلَا}

\begin{mufradat}
\kalima{قبل محرك}{أي إذا جاء بعد الميم حرف متحرك.}
\kalima{دِرَاكًا}{متتابعًا. والدال رمز \IbnKathir.}
\kalima{جَلَا}{كشف وبيّن. والجيم هنا \textbf{ليست} رمزًا لورش؛ لأنه صرّح باسم \Qalun{} قبلها.}
\end{mufradat}

\maana
إذا وقعت ميم الجمع قبل حرف متحرك فـ\IbnKathir{} يضمها ويصلها بواو في كل القرآن، و\Qalun{} مخيَّر بين الصلة والإسكان (فله وجهان).

\begin{amthila}
\ayah{أَنۡعَمۡتَ عَلَيۡهِمۡ غَيۡرِ ٱلۡمَغۡضُوبِ} \src{الفاتحة: ٧} \ila{} \IbnKathir{} (ووجه لـ\Qalun): \textbf{عليهِمُو غيرِ}.\\
\ayah{ءَأَنذَرۡتَهُمۡ أَمۡ لَمۡ تُنذِرۡهُمۡ} \src{البقرة: ٦} \ila{} \textbf{ءأنذرتهمُو أم لم تنذرهمُو لا يؤمنون}.
\end{amthila}

% ---------------------------------------------------------------------
\bayt{112}{وَمِنْ قَبْلِ هَمْزِ الْقَطْعِ صِلْهَا \qname{Nafi}{لِوَرْشِهِمْ}}{وَأَسْكَنَهَا الْبَاقُونَ بَعْدُ لِتَكْمُلَا}

\maana
\Warsh{} يصل ميم الجمع بواو \textbf{إذا كان بعدها همزة قطع} فقط (وتُمدّ عنده مدًّا منفصلًا مشبعًا). وأسكنها باقي القراء بعد \IbnKathir{} و\Qalun{} و\Warsh{}، أي: \AbuAmr{} و\IbnAmir{} و\Asim{} و\Hamza{} و\Kisai.
ونصّ الناظم على الإسكان لأن الإسكان لا يؤخذ من الضد.

\begin{qiraat}
\qira{\IbnKathir}{الصلة دائمًا قبل أي متحرك}
\qira{\Qalun}{وجهان: الصلة والإسكان قبل أي متحرك}
\qira{\Warsh}{الصلة قبل همزة القطع فقط، والإسكان قبل غيرها}
\qira{\AbuAmr، \IbnAmir، \Asim، \Hamza، \Kisai}{الإسكان دائمًا}
\end{qiraat}

\begin{amthila}
\ayah{عَلَيۡكُمۡ أَنفُسَكُمۡ} \src{المائدة: ١٠٥}، \ayah{وَمِنۡهُمۡ أُمِّيُّونَ} \src{البقرة: ٧٨} \ila{} \Warsh: \textbf{عليكمُو أنفسكم}، \textbf{ومنهمُو أميون} (مع المد المشبع).\\
أما \ayah{عَلَيۡهِمۡ غَيۡرِ} فـ\Warsh{} يسكّنها لأن بعدها غير همزة.
\end{amthila}

% ---------------------------------------------------------------------
\bayt{113}{وَمِنْ دُونِ وَصْلٍ ضَمُّهَا قَبْلَ سَاكِنٍ}{لِكُلٍّ وَبَعْدَ الْهَاءِ كَسْرُ \rz{AbuAmr}{فَتَى الْعَلَا}}

\begin{mufradat}
\kalima{فتى العلا}{\AbuAmr{} البصري، وأبوه اسمه العلاء.}
\end{mufradat}

\maana
إذا جاء بعد ميم الجمع \textbf{ساكن} فكل القراء يضمونها \textbf{بلا صلة} (لأن الصلة تُحذف لالتقاء الساكنين)، نحو: \ayah{عَلَيۡكُمُ ٱلصِّيَامُ} \src{البقرة: ١٨٣}، \ayah{وَمِنۡهُمُ ٱلَّذِينَ} \src{التوبة: ٦١}.

ثم استثنى \AbuAmr{}: فهو يكسر الميم إذا كانت بعد هاء، بشرط مذكور في البيت التالي.

\bayt{114}{مَعَ الْكَسْرِ قَبْلَ الْهَا أَوِ الْيَاءِ سَاكِنًا}{وَفِي الْوَصْلِ كَسْرُ الْهَاءِ بِالضَّمِّ \rzg{شَمْلَلَا}}

\begin{mufradat}
\kalima{شَمْلَلَا}{أسرع. والشين رمز \Hamza{} و\Kisai{} معًا.}
\end{mufradat}

\bayt{115}{كَمَا بِهِمُ الْأَسْبَابُ ثُمَّ عَلَيْهِمُ الْـ}{ـقِتَالُ وَقِفْ لِلْكُلِّ بِالْكَسْرِ مُكْمِلَا}

\begin{mufradat}
\kalima{كما}{الكاف هنا ليست رمزًا لابن عامر، بل هي كاف التشبيه وبعدها مثال (وهذا من الحروف التي «لا ريبة في اتصالها» كما قال في المقدمة).}
\kalima{مُكْمِلًا}{متمِّمًا للحكم.}
\end{mufradat}

\maana[(الأبيات ١١٣–١١٥)]
هذه حالة خاصة: ميم جمع \textbf{بعد هاء}، وقبل الهاء \textbf{كسرة أو ياء ساكنة}، وبعد الميم \textbf{ساكن}. ففيها ثلاث قراءات في الوصل:

\begin{qiraat}
\qira{\AbuAmr}{كسر الهاء وكسر الميم: \textbf{بِهِمِ الأسبابُ}، \textbf{عليهِمِ القتالُ}}
\qira{\Hamza{} و\Kisai}{ضم الهاء وضم الميم: \textbf{بِهُمُ الأسبابُ}، \textbf{عليهُمُ القتالُ}}
\qira{\Baqun}{كسر الهاء وضم الميم: \textbf{بِهِمُ الأسبابُ}، \textbf{عليهِمُ القتالُ}}
\end{qiraat}

\begin{amthila}
\begin{itemize}
\item قبل الهاء كسرة: \ayah{وَتَقَطَّعَتۡ بِهِمُ ٱلۡأَسۡبَابُ} \src{البقرة: ١٦٦}، \ayah{فِي قُلُوبِهِمُ ٱلۡعِجۡلَ} \src{البقرة: ٩٣}.
\item قبل الهاء ياء ساكنة: \ayah{كُتِبَ عَلَيۡهِمُ ٱلۡقِتَالُ} \src{البقرة: ٢٤٦}، \ayah{يُرِيهِمُ ٱللَّهُ أَعۡمَٰلَهُمۡ} \src{البقرة: ١٦٧}، \ayah{إِلَيۡهِمُ ٱثۡنَيۡنِ} \src{يس: ١٤}.
\end{itemize}
\end{amthila}

\begin{qaida}[حكم الوقف على هذه الكلمات]
عند الوقف على الميم تُسكَّن لكل القراء، وتُكسر الهاء لكل القراء (لأن سبب ضمها أو كسرها عند بعضهم هو الساكن الذي بعدها، وقد زال بالوقف): \textbf{بِهِمْ}، \textbf{يُرِيهِمْ}.
\textbf{ويُستثنى} \Hamza{} في الكلمات الثلاث \textbf{عليهم، إليهم، لديهم}؛ فإنه يضم هاءها وقفًا ووصلًا كما سبق في البيت ١١٠: \textbf{عليهُمْ}.
\end{qaida}

\begin{khulasa}[جدول جامع: ﴿عليهم القتال﴾ و﴿يريهم الله﴾]
\centering
\renewcommand{\arraystretch}{1.5}
\begin{tabular}{l c c c}
\toprule
\textbf{القارئ} & \textbf{عليهم القتال (وصلًا)} & \textbf{يريهم الله (وصلًا)} & \textbf{الوقف} \\
\midrule
\AbuAmr & عليهِمِ & يريهِمِ & عليهِمْ / يريهِمْ \\
\Hamza & عليهُمُ & يريهُمُ & عليهُمْ / يريهِمْ \\
\Kisai & عليهُمُ & يريهُمُ & عليهِمْ / يريهِمْ \\
\Baqun & عليهِمُ & يريهِمُ & عليهِمْ / يريهِمْ \\
\bottomrule
\end{tabular}
\end{khulasa}

% ---------------------------------------------------------------------
\section{خلاصة الباب}
% ---------------------------------------------------------------------
\begin{khulasa}[خلاصة سورة أم القرآن]
\renewcommand{\arraystretch}{1.45}
\begin{tabularx}{\linewidth}{>{\raggedleft\arraybackslash}p{0.2\linewidth} X}
\toprule
\textbf{الكلمة / الأصل} & \textbf{القراءات} \\
\midrule
\ayah{مَٰلِكِ} & \Asim{} و\Kisai: بالألف. \Baqun: \textbf{مَلِكِ}. \\
\ayah{ٱلصِّرَٰطَ} & \Qunbul: بالسين. \Khalaf: بالإشمام كل القرآن. \Khallad: الإشمام في الفاتحة فقط. \Baqun: بالصاد. \\
\ayah{عَلَيۡهِمۡ} وأختاها & \Hamza: بضم الهاء وقفًا ووصلًا. \Baqun: بالكسر. \\
ميم الجمع قبل متحرك & \IbnKathir: صلة. \Qalun: صلة وإسكان. \Warsh: صلة قبل الهمز فقط. \Baqun: إسكان. \\
ميم الجمع قبل ساكن & الكل بالضم بلا صلة، إلا ما بعد هاء قبلها كسر أو ياء: \AbuAmr{} يكسر الميم، و\Hamza{} و\Kisai{} يضمان الهاء والميم. \\
\bottomrule
\end{tabularx}
\end{khulasa}
               % ١٠٨–١١٥
% \include{chapters/05-idgham-kabir}       % ١١٦– ... (الباب القادم)

\backmatter
% المصادر
\chapter{المصادر}
\begin{enumerate}
\item \textbf{حرز الأماني ووجه التهاني} (متن الشاطبية)، للإمام القاسم بن فيره الشاطبي — نص الأبيات وترقيمها منه.
\item \textbf{الوافي في شرح الشاطبية}، للشيخ عبد الفتاح القاضي — المصدر الأساسي لمعاني الكلمات والشرح والأمثلة.
\item \textbf{إتحاف البرية بتحريرات الشاطبية}، للشيخ حسن خلف الحسيني — مصدر صناديق التحريرات.
\item \textbf{تقريب الشاطبية} — للمراجعة.
\end{enumerate}

\begin{tanbih}
هذا الشرح تبسيط للمصادر السابقة لغرض الدراسة. والقراءات تُؤخذ بالتلقي والمشافهة من شيخ متقن، فاعرض ما تقرؤه هنا على شيخك.
\end{tanbih}


\end{document}
